\documentclass[11pt]{report}
\usepackage[T1]{fontenc}
\usepackage{lmodern}
\usepackage[margin=1.2in]{geometry}
\usepackage{amsmath,amssymb,amsthm}
\usepackage{booktabs}
\usepackage{microtype}
\usepackage[hidelinks]{hyperref}

\newtheorem{theorem}{Theorem}[chapter]
\newtheorem{proposition}[theorem]{Proposition}
\newtheorem{lemma}[theorem]{Lemma}
\newtheorem{corollary}[theorem]{Corollary}
\theoremstyle{definition}
\newtheorem{definition}[theorem]{Definition}
\newtheorem{example}[theorem]{Example}
\theoremstyle{remark}
\newtheorem{remark}[theorem]{Remark}

\newcommand{\RR}{\mathbb{R}}
\newcommand{\EE}{\mathbb{E}}
\newcommand{\PP}{\mathbb{P}}
\newcommand{\NN}{\mathbb{N}}
\newcommand{\cA}{\mathcal{A}}
\newcommand{\cP}{\mathcal{P}}
\newcommand{\cX}{\mathcal{X}}
\newcommand{\argmax}{\operatorname*{arg\,max}}

\title{The Never-Ending Ticker\\[0.6em]\large Quantification, Attentional Jurisdiction, and the Right to Finish}
\author{Flyxion\\Independent Researcher}
\date{October 2026}

\begin{document}
\maketitle

\begin{abstract}
\noindent Contemporary information systems increasingly separate awareness from competence, measurement from value, activity from accomplishment, and responsibility from the capacity to intervene. The never-ending ticker, a continuously moving representation of a world in which the observer's actual powers remain finite, situated, and discontinuous, is taken here as the paradigmatic interface of this separation. The essay proceeds in eight parts. The first two establish that representational motion is logically independent of state change and trace the passage from measurement to incentive to behavior, with models of proxy optimization, temporal decay, and queueing. The third and fourth develop a theory of responsibility without jurisdiction and propose the Mind Your Own Business Protocol, a defeasible presumption that attention and intervention be allocated according to observational access, causal competence, feedback reliability, and legitimacy. The fifth examines institutions that are rewarded for activity and not for correction, and the sixth considers an information architecture that routes responsibility instead of broadcasting problems. The seventh part supplies philosophical foundations and a sustained treatment of the strongest objections, and the eighth proposes the right to finish as a counterprinciple to the ticker. Four formal results are proved and a fifth is stated, each with its assumptions made explicit. The essay does not argue for ignorance, for political disengagement, or against large-scale coordination. It argues that effective coordination depends on restoring the locality of practical agency and not on universalizing individual attention.
\end{abstract}

\tableofcontents

\chapter*{Conventions and Scope}
\addcontentsline{toc}{chapter}{Conventions and Scope}

The mathematical material in this essay is conceptual. Each model is introduced in order to make a particular distinction exact, to display the assumptions under which a conclusion follows, and to supply counterexamples when those assumptions fail. None of the models is offered as an empirically established law of behavior, and none of the propositions establishes an empirical conclusion by definition alone. Where the argument touches psychology or sociology, effects are presented as hypotheses requiring evidence and not as universal consequences. Quantities that are said to be normalized are assumed to be comparable by some defensible procedure, and the difficulty of supplying such a procedure is itself treated as part of the subject matter in Chapters 3 and 9.

The scholarly sources named in the text are used for the positions they are generally understood to hold, and no passage is quoted. The discussion of Jaron Lanier's argument in particular treats his image of a continuously visible running tally of individual contribution as an occasion for analysis and does not depend on the exact wording of any passage. The bibliography at the end is organized around the argument and is not offered as a survey of any of the literatures it touches.

Formal results are numbered within chapters. Four are collected for reference: the Representational Motion Theorem (Chapter 1), the Metric-Induced Goal Displacement Proposition (Chapter 3), the Attentional Saturation Proposition (Chapter 5), and the Local Corrective Allocation Proposition (Chapter 9). A fifth result, comparing broadcast and routed information networks, appears in Chapter 13 and is stated with weaker claims, since its modeling assumptions are the least settled.

\part{The Ticker as a Technological and Philosophical Object}

\chapter{The Never-Ending Ticker: An Encounter with Perpetual Motion}

\section{The Original Observation}

The investigation begins with an encounter that is deliberately trivial. In a contemporary interface for managing automated research and writing tasks, two horizontal rows of task descriptions drift across the screen in opposite directions, each row moving continuously and without pause. The tasks themselves are finite in number and fixed in content. Between one moment and the next, no task is added, completed, reordered, or altered. Nothing in the underlying system changes, and yet the display is never at rest. A viewer who glances at the screen at any arbitrary instant sees a different arrangement from the one seen an instant earlier, and the impression conveyed is of a restless, continuing, and possibly inexhaustible body of work.

The device is not mysterious, and its designers presumably adopted it for reasons that are easy to reconstruct. A scrolling row allows a larger number of items to occupy a small area than a static list would, it attracts the eye in a way a still arrangement does not, and it signals that the system is operating and has not stalled. These reasons are legitimate. What is philosophically significant is that the mechanism accomplishes its communicative effect by a purely presentational means. The signal that the system is alive is carried by the movement of the display, and the movement of the display is generated by a clock and not by the state of the world. The interface creates an impression of uninterrupted activity without any activity that it reports.

This establishes the first distinction on which the entire essay depends. Representational motion is not material progress, and the two can vary independently:
\[
\text{Representational motion} \neq \text{material progress}.
\]
The distinction is elementary, and it would hardly deserve statement if it were confined to a single decorative widget. Its interest lies in the breadth of the settings across which it recurs. Productivity software displays animated progress indicators whose motion is unrelated to the fraction of work remaining. Financial markets display price streams in which most of the changes are of a magnitude and frequency that no human decision could respond to. Social media feeds refresh with content selected to produce the impression of a world that is continuously saying something. News systems maintain banners of rolling text whose function is partly informational and partly a signal that the organization is attending to events. Institutional dashboards show counters that increment as cases are processed, and automated research environments display streams of tool calls and intermediate outputs that convey diligence while remaining silent about whether the underlying question has advanced. In each setting the ticker, understood as a continuously moving representation, occupies the visual and cognitive position that a stable record of completed and outstanding obligations might otherwise occupy.

\section{Three Meanings of the Ticker}

The word ticker names three distinct things, and the argument requires keeping them apart. The first is a mechanical device. The stock ticker of the nineteenth and early twentieth centuries was a telegraphic printer that reported changes in securities prices on a paper tape, and its name derives from the sound of its operation. Its function was to transmit, over distance and with minimal delay, a quantity that changed in the world and mattered to people who could act on it. Every character on the tape corresponded to an event at an exchange. The tape moved because prices moved.

The second meaning is an informational arrangement. The contemporary digital ticker reports changes in arbitrarily selected metrics. Because the cost of adding a metric to a display is negligible, the selection of what to display is governed less by the importance of the quantity than by its availability. A system that can count something will tend to count it, and a system that has counted something will tend to show the count. The ticker in this sense is the general form of a surface on which any measurable quantity may be placed, and its characteristic feature is that the set of quantities it reports is indifferent to whether any recipient can do anything with them.

The third meaning is a behavioral regulator. The generalized social ticker makes individuals continuously aware of what can be measured, compared, and optimized. This is the sense in which the term will most often be used in the later chapters. A ticker of this kind is a feature of the environment in which people decide what to do, and the question it raises is no longer whether it reports accurately but what it does to the process of choosing.

The crucial transition occurs when the ticker ceases to report what people are doing and begins to influence what they choose to do. The historical stock ticker could be said to have made this transition once traders began to respond to its tape in ways that moved the prices it reported, and the contemporary case differs in degree and in reach. When a displayed quantity is visible, comparable, and attached to consequences, the display enters the causal loop that produces the quantity. The device that was introduced to describe the system has become one of its components. Chapter 3 develops this transition formally. For the present it is sufficient to record that the three meanings are connected by a path of increasing causal involvement, from transmission through selection to regulation, and that the philosophical interest of the ticker lies mainly at the third stage.

\section{The Paradox of Static Information in Perpetual Motion}

The first observation can be given a precise form by separating the frequency of visual updates from the frequency of changes in the underlying system. Let $X(t)$ denote the state of an underlying system at time $t$, taking values in a set $\cX$, and let $D(t)$ denote the displayed representation of that state at time $t$, taking values in a set $\mathcal{D}$ of display configurations. A display may be said to exhibit positive presentational activity when its configuration changes with time, which in a differentiable setting means
\[
\left\|\frac{dD}{dt}\right\|>0,
\]
and the underlying system is at rest when
\[
\frac{dX}{dt}=0.
\]
The interface of the opening example satisfies both conditions at once. The display moves while the system is constant.

It is important to specify what is included in $D$. The representation includes spatial presentation as well as informational content, and not merely the symbols displayed. The displayed coordinates of otherwise unchanged objects may vary continuously, and in the example they do. A reader who confined attention to the textual content of the rows would observe that every item present at one moment is present at the next, and a reader who attended to the geometry would observe incessant change. Both observations are correct, and the paradox is the apparent coexistence of constancy and perpetual motion in a single object.

The paradox dissolves once it is recognized that the display is a function of two arguments, one of which is the state of the system and the other of which is an exogenous clock. Write the display as $D(t)=\Phi(X(t),t)$. Informational content passes through the first argument. Motion that is a function of the second argument alone carries no information about $X$, since it is generated by a process that is independent of $X$. A precise statement is the following, in which the encoding $\iota$ determines the symbols shown and the path $\pi$ determines their positions.

\begin{theorem}[Representational Motion Theorem]\label{thm:motion}
Let $\cX$ be a finite set, let $X_0$ be a random variable on $\cX$, and let the state be constant, $X(t)=X_0$ for all $t\ge0$. Let $\iota:\cX\to\mathcal{I}$ be an encoding, and let $\pi:[0,\infty)\to\RR^2$ be a continuously differentiable function that is periodic with period $T>0$, nonconstant, and known to the observer. Define the display $D(t)=(\iota(X_0),\pi(t))$. Then the following hold.
(i) The presentational path length $L(\tau)=\int_0^\tau\|\pi'(t)\|\,dt$ satisfies $L(\tau)\ge\lfloor\tau/T\rfloor\,L(T)$ with $L(T)>0$, so that presentational activity grows without bound as the observation interval lengthens.
(ii) For every $\tau\ge0$, the mutual information between the state and the entire displayed trajectory $D_{[0,\tau]}$ equals the mutual information between the state and the single initial display $D(0)$:
\[
I\bigl(X_0;D_{[0,\tau]}\bigr)=I\bigl(X_0;D(0)\bigr)=H\bigl(\iota(X_0)\bigr).
\]
Equivalently, the conditional mutual information $I\bigl(X_0;D_{(0,\tau]}\mid D(0)\bigr)$ is zero.
\end{theorem}

\begin{proof}
For (i), periodicity of $\pi$ implies periodicity of $\|\pi'\|$, so the integral over each period equals $L(T)$ and $L(\tau)\ge\lfloor\tau/T\rfloor L(T)$ because the integrand is nonnegative. Since $\pi$ is continuously differentiable and nonconstant on a period, $\pi'$ is not identically zero there, and by continuity $\|\pi'\|>0$ on an open interval, so $L(T)>0$.

For (ii), the trajectory $D_{[0,\tau]}$ is the function $t\mapsto(\iota(X_0),\pi(t))$ on $[0,\tau]$. Because $\pi$ is a fixed function known to the observer, the map that sends the value $\iota(X_0)$ to this trajectory is injective, and the trajectory determines $\iota(X_0)$ by evaluation at time $0$. Therefore $D_{[0,\tau]}$ and $D(0)$ are in bijection given the known $\pi$ and carry the same information about $X_0$. Since both are functions of $\iota(X_0)$ and each determines it, $I(X_0;D_{[0,\tau]})=I(X_0;\iota(X_0))=H(\iota(X_0))$, the final equality holding because $\iota(X_0)$ is a function of $X_0$. The conditional mutual information vanishes by the chain rule, $I(X_0;D_{[0,\tau]})=I(X_0;D(0))+I(X_0;D_{(0,\tau]}\mid D(0))$.
\end{proof}

The theorem is deliberately modest, and its modesty is the point. It establishes, under explicit assumptions, what the opening observation asserts informally: strictly positive and indefinitely accumulating presentational activity is compatible with zero information gain about the underlying state beyond that already available at the first glance. The assumptions can be relaxed in two directions. If the path $\pi$ is unknown to the observer, motion may carry information about $\pi$, though still none about $X_0$. If $\pi$ depends on $X_0$, for example by encoding a priority in the speed of a row, then motion does carry information, and the theorem does not apply. The case of interest is the former, and the proposition that the essay takes as its first formal claim is the following restatement.

\begin{proposition}[Representational Activity Without State Change]\label{prop:rep}
A display can maintain strictly positive presentational activity while conveying zero new information about its underlying state. In particular, a simple cyclic animation of an unchanging list satisfies this description.
\end{proposition}

\begin{proof}
Theorem~\ref{thm:motion} with $\pi$ a circular path, $\pi(t)=(\cos(2\pi t/T),\sin(2\pi t/T))$, gives $\|\pi'(t)\|=2\pi/T>0$ and $I(X_0;D_{(0,\tau]}\mid D(0))=0$.
\end{proof}

The relation between deterministic motion and informational novelty deserves a final emphasis. A deterministic function of time is, by definition, perfectly predictable. An observer who knows the function can compute the display at any later instant in advance, and a quantity that can be computed in advance is not news. Novelty in the technical sense requires uncertainty to be resolved, and the uncertainty about a constant state is resolved entirely at the first glance. The perpetual movement of the display therefore resolves nothing. What it does is to hold the viewer's eye, and holding the eye is an effect on the viewer, not an item of information about the world.

\section{Why This Matters}

A moving interface can suggest urgency, abundance, novelty, and incompletion without supplying evidence of any of them. Urgency is suggested because motion in the periphery of vision is a salient cue. Abundance is suggested because a stream implies a supply that is being drawn from. Novelty is suggested because a changing display resembles the display of a changing world. Incompletion is suggested because a process that continues has evidently not finished. Each of these suggestions is a statement about the underlying system, and in the opening example each is false, since the system is constant, the supply is finite, nothing is new, and the list of tasks may be entirely complete. The display does not assert these claims, and it would be uncharitable to say it lies. It insinuates them, by exploiting the ordinary association between motion and change that humans bring to the visual world.

The chapter closes with the question that governs the following seventeen. Why do modern institutions increasingly represent themselves through continuously changing displays and not through stable records of completed and outstanding obligations? A stable record has several properties that a moving display lacks. It can be compared against an earlier version of itself. It can be audited. It admits of closure, since a list from which every item has been struck is visibly finished. It can be left, because a record that is not moving does not demand to be watched. These are properties that favor deliberation, accountability, and rest. A moving display favors none of them, but it favors something else, which is the continued attention of the viewer, and a system that is financed by attention or that justifies itself by being seen to be busy has reasons to prefer it. The question is therefore not merely aesthetic. It concerns the incentives under which interfaces are designed and the philosophy of agency that those incentives silently encode.

\chapter{Jaron Lanier and the Quantification of Existence}

\section{Reading \emph{Who Owns the Future?}}

Jaron Lanier's \emph{Who Owns the Future?} \cite{lanier} provides the central historical and conceptual reference for the argument of this essay, not because it anticipates the interface of Chapter 1 but because it contains an unusually clear statement of a dilemma that interfaces of this kind embody. The book is a sustained examination of the economic arrangements of the networked world. Its principal complaint is that digital institutions extract value from the contributions of enormous numbers of people, in the form of data, writing, images, translations, and behavior, while compensating very few of them, so that the apparent abundance of free services conceals a transfer of value and a hollowing of the middle class. Its constructive proposal is an alternative economy in which the provenance of contributions is preserved and in which people are paid, perhaps in very small amounts, whenever data that originated with them is used. In the discussion of such arrangements the book considers the psychological texture of an economy in which the value of an individual's contribution is continuously and visibly accounted, and it is in this context that the image of a little ticker, a running tally of what one is worth at every moment, makes its appearance.

It would be a mistake to read Lanier as simply opposed to measurement. His economic alternative depends upon measurement, and indeed upon more measurement than existing systems perform. The wrong of the existing arrangement, in his account, is precisely that contributions go unaccounted and uncompensated. A system that recorded and rewarded them would be a more accurate and more just system. The critique is directed at an economy that counts too little in the places where counting would do justice, and the worry about the ticker arises as a worry about the consequences of repairing that deficiency.

The resulting contradiction is central to the present essay. A more equitable accounting system can also create more opportunities for compulsive optimization. If every contribution is recorded and compensated, then every contribution becomes an occasion for attending to its record, and the record becomes a stimulus. The person who is paid by the use of their data has an interest in making data that is used, and the interest gives rise to a practice of producing for the meter. The same mechanism that remedies invisible labor also renders labor continuously visible, and visibility to oneself is not a neutral condition. The tension is not peculiar to Lanier's proposal, and it recurs wherever recognition is attached to a measure that is displayed to those it measures. It is a tension and not a refutation, because the alternative of leaving contributions unaccounted is also unacceptable, and the essay's task is to ask what institutional design can preserve the benefits of accounting while declining its characteristic costs.

\section{Klout and the Gamification of Reputation}

The service known as Klout offers a compact case study. Klout was a commercial platform that computed a numerical score of an individual's social influence from activity on connected social networks, presenting the result on a bounded scale and updating it as new activity was observed. Users could view their scores, compare them with those of others, and receive notice when the score changed. The service was discontinued in 2018, and its brief history is instructive precisely because it exhibited, in a pure form, a pattern that other systems exhibit in diluted form.

The operating logic of such a score can be described in terms of what it rewards. The score aggregated signals of engagement, such as the volume of one's own posts, the responses they drew, and the standing of those who responded. A rational participant who valued the score therefore had reason to post more frequently, to seek responses, and to cultivate the attention of highly ranked others. The score was thus a measure of a person's influence and also a set of instructions for increasing it, and the second function gradually superseded the first. Observers of the service reported, and it was widely remarked, that people altered their behavior in response to it, and this is exactly the effect that the next chapter's models predict for a visible metric attached to perceived advantage.

The distinction to which the case draws attention is between an accumulating economic record and a reputation score that requires continued activity to maintain its value. The first records that something was done. The second estimates how recently and how vigorously something is being done. A person can cease contributing without losing the historical fact of their contribution, since the fact that an essay was written or a favor was done is not erased by subsequent silence. A dynamically recalculated score, by contrast, may diminish when participation stops, since it is computed from a window of recent activity. The latter introduces an implicit obligation to participate. This obligation is not imposed by any rule. It is imposed by the arithmetic of the score, and it falls upon the participant who wishes to retain what the score confers. Chapter 4 formalizes the difference as the difference between a cumulative sum and a geometrically decaying one, and draws the normative conclusion that a well-designed institution should make nonparticipation costless to recognition already earned.

\section{The Ticker as a Behavioral Environment}

The argument requires a distinction between a measurement that can be consulted and a measurement that persistently solicits attention. A bank balance may be consulted when one wishes to know it, and a published citation count may be looked up. Each of these can be transformed into a solicitation by presentation, as when a banking application sends a notification at every transaction or a platform for scholars emails its members every time a paper is cited. The numbers are identical. What differs is the relation between the number and the person, which in the first case is initiated by the person and in the second is initiated by the system.

It follows that the influence of a measure on behavior is not a function of the measure alone. A fitness score, a productivity metric, or a citation count does not exert a uniform psychological effect across all contexts. Its influence depends partly upon how it is presented, partly upon how frequently it changes, and partly upon what consequences are attached to it. A metric that changes slowly, is consulted rarely, and carries no consequence beyond information is in effect a record. A metric that changes rapidly, is pushed to the user, and is attached to rewards or sanctions is a ticker. The same underlying quantity can be either, and design choices determine which. This is a hopeful observation, since it implies that the pathologies of measurement are in important part pathologies of presentation and incentive and are therefore open to redesign.

It is also a reason for caution about causal claims. Statements to the effect that continuous metrics cause compulsive behavior generalize from observations whose scope is limited. They are plausible hypotheses, supported by an extensive body of qualitative testimony and by some experimental work on reinforcement and attention, and they require empirical examination in each context. The essay treats the behavioral effects of the ticker as hypotheses of this kind throughout.

\section{The Paradox of Equitable Measurement}

The principal question of the chapter can now be stated. Can a society recognize and compensate previously invisible contributions without making the pursuit of measurable recognition the dominant organizing principle of individual life? The question is not rhetorical, and the essay does not assume a negative answer. A positive answer would require that recognition be separable from salience, and the separation can be analyzed by distinguishing three variables that are ordinarily run together.

The first is economic recognition, the allocation of material resources and standing in proportion to contribution. The second is psychological salience, the degree to which a recognition signal occupies the attention of the person recognized. The third is institutional authority, the extent to which the recognition signal is used by others to determine opportunities, access, or treatment. A system may recognize economically without demanding attention, as a pension does that accrues without being watched. A system may be salient without conferring authority, as a game score is that has no consequence beyond play. A system may confer authority through a record that its subjects never see, as a credit file does. The three can be varied independently, and the dangerous configuration is the one in which a signal that is highly salient, continuously updated, and widely used by others to determine opportunity is attached to a person's livelihood. In that configuration a person has strong reasons to attend, to optimize, and to remain active, and the cost of declining to do so is the loss of livelihood.

The design question is therefore whether recognition can be delivered in configurations that keep salience low and authority bounded. The remainder of the essay returns repeatedly to this question, which reappears in Chapter 4 as the right to remain inactive, in Chapter 9 as the problem of the unmeasured, in Chapter 12 as the danger of a universal reputation score, and in Chapter 18 as the right to finish.

\section{Historical Precursors}

Lanier's dilemma has precursors, and it is useful to situate it briefly alongside three classical analyses of the social effects of calculation. These authors should not be presented as having anticipated digital interfaces literally. Their relevance concerns the broader tendency for formal systems of calculation to reorganize the purposes they were introduced to serve.

Georg Simmel's \emph{The Philosophy of Money} \cite{simmel} analyses money as a medium that renders heterogeneous goods commensurable, and observes that the extension of money's reach reshapes the character of the relations it mediates, making them more abstract, more calculable, and more indifferent to qualitative difference. A score that converts diverse contributions into a single number is a descendant of this operation, and the question of what is lost in the conversion is Simmel's question. Max Weber's account of rationalization \cite{weber} describes the growth of calculable, rule-governed procedures in economic and administrative life and the possibility that means designed to serve ends come to displace them, so that the procedure becomes the point. The metric that comes to define the objective is an instance of the displacement that Weber described under the heading of formal rationality. Jacques Ellul's \emph{The Technological Society} \cite{ellul} describes technique as the totality of methods rationally arrived at and having absolute efficiency in every field of human activity, and argues that the pursuit of efficiency tends to subordinate other values to itself. The ticker, as an interface optimized for the capture of attention, is a specimen of what Ellul described, though the specimen illustrates the thesis and does not confirm it.

The three citations serve a limited purpose. They show that the worry about calculation reorganizing its own purposes is old and has been formulated in more than one tradition, and they caution against supposing that the problem is a novel product of digital media. They do not show that calculation is always pernicious. Weber himself regarded rationalization as an achievement as well as a hazard, and the essay adopts that ambivalence.

\part{The Mathematics and Psychology of Measurement}

\chapter{When a Measure Becomes an Objective}

\section{Goodhart, Campbell, and Proxy Optimization}

Two aphorisms are conventionally invoked whenever measurement is used for control, and it is worth distinguishing them before using either. Charles Goodhart's observation, made in the context of British monetary policy \cite{goodhart}, was that a statistical regularity tends to collapse once pressure is placed upon it for control purposes. The usual paraphrase, that when a measure becomes a target it ceases to be a good measure, is due to Marilyn Strathern \cite{strathern} and compresses the original claim into a more general form. Donald Campbell's formulation, developed in the study of social program evaluation \cite{campbell}, is that the more a quantitative social indicator is used for social decision-making, the more subject it becomes to corruption pressures and the more apt it is to distort and corrupt the social processes it is intended to monitor. The first is primarily a statement about the deterioration of a statistical relationship. The second is primarily a statement about the behavior of people who know that an indicator is being used.

The distinction is worth drawing precisely, because the two mechanisms have different remedies. The statistical mechanism can occur in a system of passive objects. If a proxy is correlated with a target in a population, then selecting units on the proxy produces a selected group in which the correlation is weaker, and the proxy overstates the target for the selected units. The behavioral mechanism requires agents who respond to the measurement. If the proxy can be raised by actions that do not raise the target, then agents who benefit from a high proxy will take those actions, and the proxy will cease to correlate with the target in the population that results. The first deteriorates because selection is imperfect, and the second deteriorates because the measured have an interest in the measure.

The statistical mechanism can be exhibited exactly in a standard Gaussian model.

\begin{proposition}[Regressional Overstatement]\label{prop:regress}
Let $V\sim\mathcal N(0,\sigma_V^2)$ be the underlying value of a unit, let $\varepsilon\sim\mathcal N(0,\sigma_\varepsilon^2)$ be independent noise, and let $U=V+\varepsilon$ be the measured proxy. Then
\[
\EE[V\mid U=u]=r\,u,\qquad r=\frac{\sigma_V^2}{\sigma_V^2+\sigma_\varepsilon^2}\in(0,1)\quad\text{whenever }\sigma_\varepsilon>0 .
\]
Consequently any unit selected because its proxy value is $u>0$ has expected true value $ru<u$, and the overstatement $u-\EE[V\mid U=u]=(1-r)u$ grows linearly in the selection threshold.
\end{proposition}

\begin{proof}
The pair $(V,U)$ is jointly Gaussian with $\mathrm{Cov}(V,U)=\sigma_V^2$ and $\mathrm{Var}(U)=\sigma_V^2+\sigma_\varepsilon^2$. The conditional expectation of one component of a bivariate Gaussian given the other is linear, with slope $\mathrm{Cov}(V,U)/\mathrm{Var}(U)=r$ and, since both means are zero, zero intercept.
\end{proof}

The proposition shows that even without any strategic behavior, optimization against a noisy proxy produces disappointment proportional to the strength of the optimization. It supplies the statistical reading of Goodhart's law and nothing more. The behavioral reading is addressed in Section 3.4, where agents control the proxy.

\section{Value and Measurement}

Let $V(a)$ represent the underlying value of an activity $a$, and let $M(a)$ represent its measured value, for $a$ in a finite set $\cA$ of available activities. Define a stylized decision rule by which a decision-maker, at time $t$, selects
\[
a_t^\ast=\argmax_{a\in\cA}\left[(1-\alpha_t)\,V(a)+\alpha_t\,M(a)\right],
\]
where $\alpha_t\in[0,1]$ represents the influence of the visible measurement upon decision-making. At $\alpha_t=0$ the decision-maker pursues value alone, and at $\alpha_t=1$ the decision-maker pursues the measure alone. The expression requires comparable normalization of $V$ and $M$, since a weighted sum of quantities in incommensurable units is meaningless, and it is a conceptual model and not an empirically established law of behavior. Its function is to isolate a single variable, the weight on the measure, and to ask what follows from changes in that variable.

Two results are immediate and clarify the structure of the problem. The first concerns the threshold at which a decision-maker switches from one activity to another. Suppose that two activities $a_1$ and $a_2$ satisfy $V(a_1)>V(a_2)$ and $M(a_2)>M(a_1)$, so that the first is worth more and the second measures better. Setting the two weighted sums equal gives $(1-\alpha)(V_1-V_2)=\alpha(M_2-M_1)$, and therefore the switching weight
\[
\alpha^\ast=\frac{V_1-V_2}{(V_1-V_2)+(M_2-M_1)} .
\]
For $\alpha<\alpha^\ast$ the valuable activity is chosen, and for $\alpha>\alpha^\ast$ the better-measuring one is. The threshold is small when the value difference is small relative to the measurement difference, which is to say when an activity that is only slightly less valuable is substantially more visible. The practical content of this is that small differences in value are easily overwhelmed by modest weight on the measure.

The second result is general and concerns the direction of change.

\begin{proposition}[Monotone Displacement]\label{prop:monotone}
Let $\alpha<\alpha'$ and let $a$ and $a'$ be maximizers of the weighted objective at weights $\alpha$ and $\alpha'$ respectively. Then $V(a')\le V(a)$ and $M(a')\ge M(a)$.
\end{proposition}

\begin{proof}
Write $\Delta V=V(a')-V(a)$ and $\Delta M=M(a')-M(a)$, and define $f(\beta)=(1-\beta)\Delta V+\beta\Delta M$, a linear function of $\beta$. Optimality of $a$ at $\alpha$ gives $f(\alpha)\le0$, and optimality of $a'$ at $\alpha'$ gives $f(\alpha')\ge0$. A linear function that is nonpositive at $\alpha$ and nonnegative at $\alpha'>\alpha$ is nondecreasing, so $f(0)\le f(\alpha)\le0$ and $f(1)\ge f(\alpha')\ge0$. Since $f(0)=\Delta V$ and $f(1)=\Delta M$, the claims follow.
\end{proof}

The proposition says that as the measure gains influence, the chosen activity never improves in true value and never worsens in measured value. This is the structure of displacement in its barest form, and it implies that the effect of a visible measure on a decision-maker is directional, regardless of how the specific activities are arranged.

The substantive question is under what conditions $\alpha_t$ increases. Five mechanisms deserve examination. Reward raises $\alpha$ when the measure is attached to material or social benefit, since a decision-maker who is paid by the measure has a reason to weight it. Repetition raises $\alpha$ when exposure habituates attention to the measure and renders its changes salient. Competition raises $\alpha$ when the relevant comparison is with others, because a relative measure punishes the one who ignores it even if the absolute value of ignoring it is small. Habituation raises $\alpha$ when what began as an instrument becomes a standing concern. Institutional necessity raises $\alpha$ when an organization requires the measure for access, funding, or standing, so that the decision-maker must weight it whatever her private view of its relation to value. Each of these is an empirical hypothesis about particular settings, and the model does not decide among them.

\section{Endogenous Preferences}

The interesting possibility is not simply that people optimize the wrong variable. Repeated exposure to the measurement may change what they experience as desirable. This is a more troubling hypothesis than strategic misdirection, because strategic misdirection can be reversed by removing the incentive, whereas a changed preference may persist after the incentive has gone. The chapter should distinguish deliberate strategic optimization, in which the decision-maker knows that she is serving a proxy, from the gradual internalization of institutional metrics, in which the metric becomes part of what the decision-maker wants.

A minimal model captures the second possibility. Let the weight on the measure evolve according to
\[
\alpha_{t+1}=\alpha_t+\lambda\,s_t\,(1-\alpha_t),\qquad \alpha_0\in[0,1),\quad \lambda\in(0,1],\quad s_t\in[0,1],
\]
where $s_t$ is the intensity of exposure to the measure at time $t$ and $\lambda$ is a susceptibility parameter. The model asserts that exposure moves the weight toward its maximum at a rate proportional to the remaining distance.

\begin{proposition}[Internalization]\label{prop:internal}
In the preceding model, $1-\alpha_t=(1-\alpha_0)\prod_{k<t}(1-\lambda s_k)$. If $\sum_k s_k=\infty$ then $\alpha_t\to1$. If $\sum_k s_k<\infty$ then $\alpha_t$ converges to a limit strictly less than $1$, provided $\lambda s_k<1$ for all $k$.
\end{proposition}

\begin{proof}
The product formula follows by induction from $1-\alpha_{t+1}=(1-\alpha_t)(1-\lambda s_t)$. For factors $x_k=\lambda s_k\in[0,1)$, the infinite product $\prod(1-x_k)$ converges to zero if and only if $\sum x_k=\infty$, and converges to a strictly positive limit if and only if $\sum x_k<\infty$. Both statements are standard, and follow from the inequalities $1-x\le e^{-x}$ and $1-x\ge e^{-x-x^2}$ for $x\in[0,1/2]$, together with the fact that only finitely many $x_k$ can exceed $1/2$ when their sum converges.
\end{proof}

The proposition has an interpretive moral that is independent of the particular dynamics assumed. If exposure is sustained at any positive level, then the weight on the measure approaches its maximum, and the approach occurs however small the susceptibility. Only exposure that is summable, which is to say exposure that eventually ceases or becomes vanishingly rare, leaves a residue of weight on value. A continuous ticker is by construction an exposure that does not cease. The model does not prove that people internalize metrics, and it should not be mistaken for evidence. It shows what internalization would imply if the mechanism of exposure-driven drift exists, and it identifies the design variable that matters, namely the integral of exposure.

\section{Measurement-Induced Behavioral Distortion}

The feedback structure of the problem can be written as a system of three maps. The measured value is a function of the state, the behavior is updated in response to the measure, and the state is changed by behavior:
\[
M_t=f(X_t),\qquad A_{t+1}=g(A_t,M_t),\qquad X_{t+1}=h(X_t,A_{t+1}).
\]
The observed metric influences behavior, and behavior subsequently changes the measured state. It is the closing of this loop that distinguishes a measuring instrument from a regulator. Examination of fixed points, of the amplification of perturbations, and of the possibility of self-reinforcing but socially unproductive equilibria proceeds by linearization.

Consider the scalar linear case in which $f$ is the identity, so that $M_t=X_t$, and let
\[
A_{t+1}=\rho A_t+kM_t,\qquad X_{t+1}=X_t+bA_{t+1},
\]
with $\rho\in[0,1)$ capturing the persistence of behavior, $k$ the responsiveness of behavior to the metric, and $b$ the effect of behavior on the state. Substituting gives the system on the vector $(X_t,A_t)$ with matrix
\[
\begin{pmatrix}1+bk & b\rho\\ k&\rho\end{pmatrix},
\]
whose characteristic polynomial is $p(\lambda)=\lambda^2-(1+bk+\rho)\lambda+\rho$, since the determinant is $(1+bk)\rho-bk\rho=\rho$ and the trace is $1+bk+\rho$.

\begin{proposition}[Amplifying Feedback]\label{prop:amplify}
If $bk>0$, so that behavior responds positively to the metric and positively affects it, then $p(1)=-bk<0$, and therefore the matrix has a real eigenvalue exceeding $1$. For generic initial conditions the measured state diverges. If $bk<0$, so that the response is corrective, then $p(1)=-bk>0$ and the unit eigenvalue is not a root, and stability depends on the remaining conditions $p(-1)>0$ and $\rho<1$.
\end{proposition}

\begin{proof}
Evaluating, $p(1)=1-(1+bk+\rho)+\rho=-bk$. Since $p$ is a monic quadratic, $p(\lambda)\to+\infty$ as $\lambda\to\infty$, so if $p(1)<0$ there is a root greater than $1$. The corresponding eigenvector has nonzero projection on generic initial data, whose trajectory therefore grows geometrically in the direction of that eigenvector. The case $bk<0$ is the standard Jury criterion: both roots of a monic quadratic lie inside the unit circle if and only if $p(1)>0$, $p(-1)>0$, and $|p(0)|<1$; here $p(0)=\rho<1$, and $p(-1)=2+bk+2\rho$.
\end{proof}

The result is again conditional and should be read as a pattern and not as a prediction. In the amplifying case the metric grows without bound, which in a real system is tempered by saturation, and the interesting equilibria are those at which growth is stopped by a ceiling or exhausted by costs. What the result does show is that a positive loop between a visible metric and the behavior that raises it has the origin as its only fixed point, which is unstable, and that any stable configuration is therefore a boundary configuration: the metric is maxed out, or the behavior is exhausted. Whether such a configuration is socially unproductive depends on whether the metric tracks anything of value at the boundary, and that dependence is the subject of the next result.

\subsection*{Metric-Induced Goal Displacement}

The behavioral reading of Goodhart and Campbell requires a model in which an agent allocates effort between a task that is measured and a task that is not. The following multitask model is a simplified instance of the principal-agent analyses of Holmstr\"om and Milgrom \cite{holmstrom}, and is closely related to Kerr's classic observation about rewarding one behavior while hoping for another \cite{kerr}.

An agent chooses nonnegative efforts $e_1$ on a measured task and $e_2$ on an unmeasured task. The agent's cost of effort is
\[
C(e_1,e_2)=\tfrac12\left(e_1^2+e_2^2\right)+\delta\,e_1e_2,\qquad \delta\in(0,1),
\]
where $\delta$ represents substitutability: effort on one task raises the marginal cost of effort on the other. The agent receives a reward $\beta_1 e_1$ proportional to the measured task, and has an intrinsic motivation of strength $\eta>0$ for the unmeasured task, so that the agent's utility is $\beta_1e_1+\eta e_2-C(e_1,e_2)$. True value is $W=v_1e_1+v_2e_2$ with $v_1,v_2>0$, which counts the value of output and abstracts from the cost of effort. The measured performance is $M=e_1$.

\begin{proposition}[Metric-Induced Goal Displacement]\label{prop:displace}
Let $\beta_1$ increase from zero. The agent's optimal efforts are as follows.
For $\beta_1\le\delta\eta$: $e_1=0$ and $e_2=\eta$.
For $\delta\eta<\beta_1<\eta/\delta$: both efforts are interior, with
\[
e_1=\frac{\beta_1-\delta\eta}{1-\delta^2},\qquad e_2=\frac{\eta-\delta\beta_1}{1-\delta^2}.
\]
For $\beta_1\ge\eta/\delta$: $e_1=\beta_1$ and $e_2=0$.
Throughout the interior range, measured performance rises at rate $1/(1-\delta^2)$ while true value changes at rate $(v_1-\delta v_2)/(1-\delta^2)$. Hence, if $v_1<\delta v_2$, then raising the reward increases measured performance and strictly lowers true value, and moreover $W(\beta_1)<W(0)$ for every $\beta_1\in(\delta\eta,\,\eta/\delta]$.
\end{proposition}

\begin{proof}
The first-order conditions for an interior optimum are $\beta_1-e_1-\delta e_2=0$ and $\eta-e_2-\delta e_1=0$. Solving the linear system, whose determinant is $1-\delta^2>0$, gives the stated interior efforts. The interior solution has $e_1>0$ if and only if $\beta_1>\delta\eta$ and $e_2>0$ if and only if $\beta_1<\eta/\delta$. When $\beta_1\le\delta\eta$ the unconstrained solution would have $e_1\le0$, so $e_1=0$ and the problem reduces to maximizing $\eta e_2-\tfrac12e_2^2$, giving $e_2=\eta$. When $\beta_1\ge\eta/\delta$ the unconstrained solution would have $e_2\le0$, so $e_2=0$ and $e_1=\beta_1$. The solution is continuous across the boundaries. On the interior range, $de_1/d\beta_1=1/(1-\delta^2)$ and $de_2/d\beta_1=-\delta/(1-\delta^2)$, so that $dW/d\beta_1=(v_1-\delta v_2)/(1-\delta^2)$. At $\beta_1=\delta\eta$ the value is $W=v_2\eta=W(0)$, and if the slope is negative then $W<W(0)$ throughout the interior. At $\beta_1=\eta/\delta$, $W=v_1\eta/\delta$, which is less than $v_2\eta$ exactly when $v_1<\delta v_2$.
\end{proof}

Two observations qualify the result. First, the displacement is not permanent in this model: for sufficiently large $\beta_1$ the unmeasured task has already been abandoned entirely, and further increases in the reward raise true value again through the measured task alone, at rate $v_1$. Welfare is thus non-monotone in the strength of the incentive, and the interval of harm is the interval of partial displacement. This corresponds to the common experience that a moderate pressure on a metric damages work that a strong enough pressure would simply replace. Second, the result concerns the incentive attached to the proxy, and does not by itself indicate what should be done. Chapter 11 returns to the question of when independent verification, which would allow a reward to be attached to the unmeasured task as well, prevents the displacement.

\section{Measurement Is Not Inherently Pathological}

The preceding analysis could be misread as an argument against measurement, and counterexamples show why the misreading would be a mistake. A clinical thermometer measures body temperature, and the fact that it can be read does not distort the physiological processes it reports. Engineering diagnostics detect fatigue cracks in structural elements, and the measurement is valuable precisely because it is tied to a corrective response that can be taken before failure. Safety-critical systems, such as the monitoring of reactor coolant, display a continuously updated quantity and are justified in doing so. In each case, the measurement improves decision-making because it reliably identifies an actionable deviation and supports a corrective response.

The common feature of these cases is not that they are unmeasured by humans, nor that they change slowly. Rather, the measurement is subordinate to the purpose for which it was introduced. The deviation that it flags corresponds to a state of affairs that the observer can do something about, there is a known standard against which deviation is judged, and the consequence of the measurement is a decision whose merits are assessed by outcomes in the world and not by the measurement. The distinguishing question is therefore not whether something is measured, but whether the measurement remains subordinate to its purpose. In the language of the decision rule above, a healthy measurement is one for which the weight $\alpha_t$ is kept low by structure and not by willpower, because the measure is consulted when a decision requires it and is not present otherwise.

\chapter{The Temporal Architecture of Compulsion}

\section{Accumulating, Decaying, and Continuously Recalculated Metrics}

Metrics differ in how they treat the past, and the difference has consequences for the obligations they impose. A metric that preserves historical accomplishment records contributions and retains them. For such an accumulating record, with contribution $c_t\ge0$ at time $t$,
\[
M_{t+1}=M_t+c_t,\qquad c_t\ge0.
\]
A metric that requires continual activity to maintain a score instead discounts the past at a geometric rate. For a decaying reputation measure with retention factor $\rho$,
\[
M_{t+1}=\rho M_t+c_t,\qquad 0<\rho<1 .
\]
Under the second arrangement, inactivity produces decline even when the individual has committed no error. The score of a person who stops contributing at time $T$ evolves as $M_{T+n}=\rho^nM_T$, which tends to zero. Such a score has a half-life of $\ln 2/\ln(1/\rho)$ periods.

Two properties distinguish the arrangements quantitatively. In the accumulating case a constant contribution rate $c$ produces linear growth without limit, and the record is an integral of effort. In the decaying case a constant contribution rate produces convergence to the steady state $M^\ast=c/(1-\rho)$, and the score measures a rate of activity and not a stock of accomplishment. It follows that holding a decaying score at a fixed level $m>0$ requires a perpetual contribution of $(1-\rho)m$ in every period. The obligation to maintain is a property of the arithmetic of the score. It is independent of whether the person has anything further to say, do, or give.

\begin{proposition}[Perpetual Maintenance]\label{prop:maintain}
In the decaying model, a score cannot be held at any positive level $m$ by a finite number of nonzero contributions. In the accumulating model, a record at level $m$ is retained indefinitely with zero further contribution.
\end{proposition}

\begin{proof}
In the decaying model, $M_{t+1}=\rho M_t+c_t$ with $c_t=0$ for all $t$ beyond some finite $T$ gives $M_t=\rho^{t-T}M_T\to0$, which is below any positive level $m$ eventually. In the accumulating model, $c_t=0$ gives $M_{t+1}=M_t$.
\end{proof}

\section{The Institutional Meaning of Decay}

Decay functions can be justified in contexts where recency matters, and it would be mistaken to treat them as inherently illegitimate. The standard justification is epistemic. Suppose that the quality $x_t$ of a person's current competence drifts over time as a random walk, $x_{t+1}=x_t+w_t$ with $w_t$ a zero-mean disturbance, and that each period yields a noisy observation $y_t=x_t+v_t$. The estimator of current competence that minimizes mean squared error in the steady state is the Kalman filter for the local-level model, which takes the form of exponential smoothing,
\[
\hat x_{t+1}=(1-K)\hat x_t+K\,y_{t+1},
\]
with gain $K\in(0,1)$ determined by the ratio of the disturbance variance to the observation noise variance. This has the same recursive form as the decaying metric with $\rho=1-K$. A system that estimates how good a person is now, in a domain where skill or relevance drifts, has reason to weight recent evidence more heavily than distant evidence, and the weighting is a feature of inference and not a disciplinary device.

The same arithmetic, however, can be used in a different capacity. When the estimate of present competence is displayed as a score to the person and attached to opportunities, the estimator becomes a demand. A person who must maintain a score to retain access has a reason to produce observations that raise the estimate, and the observations produced for that purpose are no longer the same evidence the estimator was designed to process. The estimator thereby loses the property that justified its decay. The legitimate use of decay is therefore for inference about the present state of affairs, in settings where the present state is what a decision requires, and the illegitimate use is the manufacture of perpetual obligation. The distinction can be put as a question to be asked of any decaying measure: does a decision depend on the estimate of the present, or does the system depend on the continued activity of the person estimated? Where it is the second, the decay serves the system's interest in engagement and borrows the authority of inference.

\section{Intermittent Rewards and Continuous Anticipation}

The psychology of reinforcement is sometimes invoked in discussions of digital interfaces, and the invocation requires care. Experimental work with animals, beginning with the schedules of reinforcement catalogued by Ferster and Skinner \cite{ferster}, established that rewards delivered on variable schedules, in which the number of responses required for reward varies unpredictably, tend to sustain high and persistent response rates and to be resistant to extinction. This is an experimentally established result in laboratory settings with specific species and specific tasks.

The step from this result to claims about particular digital interfaces is a speculative application and not an established finding. An interface that delivers social approval, novel content, or messages at irregular intervals shares a structural feature with a variable schedule, and it has been argued that the resemblance explains compulsive checking. The argument is plausible, and the essay treats it as a hypothesis worth testing. It remains true that the contexts differ from the laboratory in many respects, including the nature of the reward, the availability of alternatives, the presence of competing demands, and the cognitive sophistication of the participants. The essay's argument does not depend on the hypothesis, and it is stated here only so that its status may be recorded. The temporal architecture of a metric can create obligations through arithmetic alone, as the preceding sections showed, and need not draw on any special psychological mechanism.

\section{The Right to Remain Inactive}

The concept of nonparticipation without penalty can now be made precise. Ask whether a well-designed informational institution should permit users to stop attending to a metric without losing recognition for past accomplishments or forfeiting unrelated opportunities. The proposed answer is affirmative, and the principle can be defined as a formal property of recognition functions.

Consider a recognition function that is linear in past contributions, $R_T=\sum_{s<T}w_{T,s}\,c_s$, where $w_{T,s}\ge0$ is the weight at time $T$ of the contribution made at time $s$. Say that the function is \emph{inactivity-invariant} if, for every history of contributions, ceasing to contribute never lowers recognition, which is to say that $R_{T+1}\ge R_T$ whenever $c_T=0$.

\begin{proposition}[Characterization of Inactivity Invariance]\label{prop:inactive}
A linear recognition function is inactivity-invariant if and only if $w_{T+1,s}\ge w_{T,s}$ for all $s<T$. The accumulating record ($w_{T,s}=1$) is inactivity-invariant, and the decaying score ($w_{T,s}=\rho^{T-1-s}$ with $\rho<1$) is not, whenever any past contribution is positive.
\end{proposition}

\begin{proof}
With $c_T=0$, $R_{T+1}-R_T=\sum_{s<T}(w_{T+1,s}-w_{T,s})c_s$. If $w_{T+1,s}\ge w_{T,s}$ for all $s$ then the difference is nonnegative for all histories with $c_s\ge0$. Conversely, if $w_{T+1,s_0}<w_{T,s_0}$ for some $s_0$, the history that has only $c_{s_0}=1$ yields a negative difference. For the accumulating record the weights are constant, so the condition holds with equality. For the decaying score, $w_{T+1,s}=\rho\,w_{T,s}<w_{T,s}$ whenever $w_{T,s}>0$.
\end{proof}

The right to remain inactive then has an institutional implementation that does not require abandoning decay altogether. Two registers can be separated. A record of accomplishment is kept in an accumulating form, is inactivity-invariant, and is the basis of recognition and of access to opportunity that depends on what a person has done. A separate estimate of present relevance, where relevance is genuinely needed for a decision, is computed with decay but is not presented as a standing score to be maintained, and is consulted only at the moment of the decision for which it is relevant. The two registers answer two different questions, what has this person done and what is this person likely to do now, and conflating them is the source of the compulsion. This principle becomes a central normative element of the protocol developed in Part IV.

\chapter{Information Overload and the Scarcity of Attention}

\section{Herbert Simon and the Attention Economy}

Herbert Simon's observation, made in a discussion of organizational design for an information-rich world \cite{simon71}, was that a wealth of information creates a poverty of attention and a need to allocate that attention efficiently among the overabundance of information sources that might consume it. The statement is often cited as the origin of the attention economy, and it identifies a structural point that is easily lost. The availability of information is a property of the environment. The capacity to process it is a property of the agent. These can diverge without limit, since the first is largely a matter of storage and transmission costs, which have fallen dramatically, while the second is bounded by the biology of the human nervous system, which has not.

Simon's work on bounded rationality \cite{marchsimon} adds that organizations are, among other things, devices for coping with this divergence. A hierarchy, a division of labor, and a system of rules each reduce the amount of information that any one member must handle by routing, filtering, and standardizing. The significance for the present argument is that the problem of attention is not merely a problem of personal discipline. It is an organizational and institutional problem, and its solution lies partly in the design of channels. An information environment that eliminates the filters of organizations does not thereby increase the information available to a person in any useful sense, since the person's capacity does not change. It increases the volume of unfiltered material that she must either process or ignore.

\section{A Queueing Model of Incoming Information}

Let $\lambda$ denote the arrival rate of informational demands and $\mu$ the observer's effective processing capacity. For a simplified queue with suitable assumptions, persistent $\lambda\ge\mu$ prevents a stable finite backlog. The standard model is the single-server queue with Poisson arrivals and exponentially distributed service times, the M/M/1 queue \cite{kleinrock}. Its steady-state existence condition is $\rho=\lambda/\mu<1$, in which case the number in the system has a geometric distribution with mean $\rho/(1-\rho)$, and the mean time in system, by Little's law \cite{little}, is $1/(\mu-\lambda)$. When $\lambda\ge\mu$, the queue length grows without bound and no stationary distribution exists.

The result does not directly describe every social feed, and the limitation should be stated explicitly. Interfaces may discard, rank, suppress, or expire messages and not preserve them in a queue. An infinite buffer is an idealization that fits email better than it fits a feed, in which older items are simply displaced from view. The queueing model is appropriate where messages persist and impose an unprocessed obligation, as in a mailbox or a task tracker, and the discard model is appropriate where messages vanish. In the discard case the relevant question is not the length of the backlog but the probability that a given message is ever considered.

\section{The Displacement of Consequential Messages}

The structure of the problem is most clearly seen when messages differ in importance. Suppose that two classes of message arrive at the observer, independent Poisson streams of rates $\lambda_a$ for actionable messages and $\lambda_n$ for non-actionable ones, with total rate $\lambda=\lambda_a+\lambda_n$. Suppose that the observer cannot distinguish them on arrival and serves in order of arrival at exponential rate $\mu$. An actionable message has a deadline $\tau$ by which it must receive consideration to be useful. A particularly important case involves a direct request from an actual collaborator being obscured by a much larger volume of automated updates, and the model addresses exactly this.

\begin{proposition}[Attentional Saturation]\label{prop:saturation}
Under the preceding assumptions (Poisson arrivals, exponential service of rate $\mu$, first-come-first-served discipline, infinite buffer, classes indistinguishable on arrival) and with $\lambda<\mu$, the time from arrival to completion of consideration of any message is exponentially distributed with rate $\mu-\lambda$. The probability that an actionable message is considered within its deadline $\tau$ is therefore
\[
\PP(\text{timely})=1-e^{-(\mu-\lambda_a-\lambda_n)\tau},
\]
which is strictly decreasing in $\lambda_n$ and tends to $0$ as $\lambda_n\uparrow\mu-\lambda_a$. If $\lambda\ge\mu$ then the delay grows without bound and the probability of timely consideration of an arbitrary message tends to $0$ as time passes. If instead the observer can classify messages perfectly on arrival and serves actionable messages with preemptive-resume priority, then the sojourn time of an actionable message is exponential with rate $\mu-\lambda_a$ and is independent of $\lambda_n$.
\end{proposition}

\begin{proof}
For the M/M/1 queue under first-come-first-served service, the stationary sojourn time is exponential with rate $\mu-\lambda$, a standard result obtained by conditioning on the number found in the system on arrival, which is geometric by the PASTA property, and summing the resulting mixture of Erlang service times. The timely probability is then the exponential distribution function at $\tau$. Differentiation with respect to $\lambda_n$ shows strict monotonic decrease, and the limit is immediate. When $\lambda\ge\mu$ the system is unstable and the queue length diverges almost surely, so the sojourn time of the $n$th arrival grows without bound. Under preemptive-resume priority with exponential service, the high-priority class is unaffected by the low-priority class, since any service of a low-priority message is interrupted on arrival of a high-priority one and resumed without loss by memorylessness, so the high-priority class experiences an M/M/1 queue with arrival rate $\lambda_a$ alone.
\end{proof}

The proposition distinguishes sharply between two cases. If the observer cannot tell which messages matter before reading them, then every additional irrelevant message degrades the timeliness of the relevant ones, and the degradation is catastrophic as the total load approaches capacity. If the observer or her system can classify perfectly, irrelevant volume is harmless to the important class. The reality lies between, and the dependence on classification quality is the important lesson. A sufficiently accurate filter can make noise costless, an imperfect filter makes noise costly in proportion to its errors, and the absence of a filter makes it maximally costly. Improving filters is accordingly a more promising response to overload than appeals to individual discipline, and much of Part VI is concerned with what a good filter would be.

A weaker result holds for the discard model. If the observer considers $k$ of $N$ equally available messages chosen uniformly at random, then the probability that a particular message is considered is $\min(1,k/N)$, which decreases in $N$ and, for fixed $k$, vanishes as the volume of other messages grows. The conclusion is the same in structure, and it again requires that ranking be uninformative about importance. A ranking that raises important messages above others changes the result, and a ranking that raises engaging messages above others may reduce the probability for the important ones.

\section{Notification Saturation and Rational Disengagement}

The model also explains why disabling notifications may be a reasonable adaptation to informational overload and not evidence of indifference. A person who observes that most notifications are non-actionable, and who cannot determine which are actionable without reading them, faces the choice between reading all, which incurs the cost of saturation, and reading none or few, which incurs the cost of missing the rare actionable one. If the proportion of actionable messages is low enough that the cost of reading all exceeds the expected cost of missing a few, the rational policy is to read few. The action of muting is an inference about the information environment and cannot be interpreted as a verdict on the importance of the matters notified.

The institution may respond to reduced attention by increasing notification frequency. A platform whose revenue or purpose depends on engagement observes that response to its messages has declined and raises the rate of messages in order to restore engagement. This is a feedback loop, and its structure can be analyzed. Let $r(\lambda)$ be the user's response rate to messages when the rate of messages is $\lambda$, assumed decreasing, and let the platform adjust $\lambda_{t+1}=\lambda_t+\kappa\,(r^\ast-r(\lambda_t))$, with a target response rate $r^\ast$ and gain $\kappa>0$.

\begin{proposition}[Counterproductive Escalation]\label{prop:escalate}
If $r$ is strictly decreasing and $r(\lambda_0)<r^\ast$, then $\lambda_t$ is strictly increasing in $t$, no fixed point exists in $[\lambda_0,\infty)$, and the sequence either diverges or reaches the upper limit of the feasible range.
\end{proposition}

\begin{proof}
We show by induction that $\lambda_{t+1}>\lambda_t\ge\lambda_0$. Since $r$ is decreasing, $\lambda_t\ge\lambda_0$ implies $r(\lambda_t)\le r(\lambda_0)<r^\ast$, so $\lambda_{t+1}-\lambda_t=\kappa(r^\ast-r(\lambda_t))>0$. A fixed point $\lambda^\dagger\ge\lambda_0$ would require $r(\lambda^\dagger)=r^\ast$, but $r(\lambda^\dagger)\le r(\lambda_0)<r^\ast$. A strictly increasing sequence with no fixed point cannot converge to a finite limit in the interior, since by continuity of the update map a limit would be a fixed point (when $r$ is continuous), and so it is unbounded or terminates at a boundary.
\end{proof}

The proposition shows that a response-targeting platform in a setting where messages crowd out responses has no stable equilibrium short of the boundary. It is a model of what happens when an institution measures its success by attention and treats declining attention as a deficiency to be repaired by more demands, and it is silent about whether any actual platform behaves in this way.

\section{The Difference Between Being Informed and Being Capable}

The conclusion of the chapter is that awareness is not an adequate proxy for agency. An observer may know about thousands of problems while possessing the practical capacity to address only a handful, and the interface that presents the thousands does not transfer capacity to the observer in the course of transferring awareness. The unstated assumption in many designs of information systems is that increased visibility of a problem increases the probability that it will be addressed, because a visible problem has more potential helpers. The assumption holds when the helpers who see the problem are the ones capable of helping it. It fails when the number of viewers who cannot help vastly exceeds the number who can, because then the expected number of effective responders is a tiny fraction of the audience and the remainder are left with awareness and no outlet. The next part examines what happens to those people.

\part{The Crisis of Universalized Attention}

\chapter{Responsibility Without Jurisdiction}

\section{The Globalization of Concern}

Contemporary information networks expose individuals to problems across geographic, institutional, and social boundaries. A person with a telephone can, within minutes, observe a famine, a legal injustice, a corporate scandal, an environmental disaster, and the outcome of a municipal dispute in a city she has never visited. The exposure is not new in kind, since newspapers, radio, and television extended the reach of attention long before the network, but its density, immediacy, and personalization are new. The relevant change is not simply increased knowledge, but increased exposure to implied demands for response. A problem presented to a person is presented with a framing, and the framing frequently carries the suggestion that the person is positioned to do something about it, whether by donating, sharing, signing, protesting, boycotting, or at minimum by feeling the appropriate concern.

The suggestion is often false in the sense that matters. The person to whom a distant crisis is presented usually has no authority over its resolution, no access to the relevant causal levers, no means of verifying the claims made about it, and no way of observing the consequences of anything she might do. The presentation nevertheless shapes her sense of what is incumbent upon her, and the result is a systematic production of felt obligation that has no corresponding capacity. This is the condition the chapter examines.

\section{Moral Awareness and Practical Responsibility}

Four relations between an agent and a problem are frequently conflated, and the analysis requires that they be separated. Write $\mathrm{Know}(a,p)$ for the claim that agent $a$ knows that problem $p$ exists, $\mathrm{Care}(a,p)$ for the claim that $a$ has reason to care about $p$, $\mathrm{Duty}(a,p)$ for the claim that $a$ possesses a duty concerning $p$, and $\mathrm{Able}(a,p)$ for the claim that $a$ has the capacity to correct $p$. These relations frequently overlap but are not identical, and each of the sixteen combinations of truth values is instantiated by some familiar case. The following table lists representative cases for several of them.

\begin{table}[ht]
\centering
\small
\begin{tabular}{@{}cccp{6.2cm}@{}}
\toprule
Know & Care & Duty & Able and representative case\\
\midrule
yes & yes & no & no: a private citizen aware of a distant famine, with reason to care and no assigned duty and no means of relief\\
yes & yes & yes & yes: a physician aware of a patient's deterioration in her own ward\\
yes & yes & yes & no: an official aware of a deficiency whose remedy has been defunded\\
yes & no & yes & yes: a contractor aware of a defect who is indifferent but bound to repair it\\
no & yes & yes & yes: a supervisor responsible for a unit who has not been informed of a fault in it\\
no & no & no & yes: a skilled bystander who could help but has no inkling that help is needed\\
\bottomrule
\end{tabular}
\caption{Illustrative independence of knowledge, care, duty, and ability.}
\label{tab:fourfold}
\end{table}

The independence of the four is a logical point with practical consequences. From knowledge it does not follow that one has a reason to care, since one can know of matters that are trivial. From reason to care it does not follow that one has a duty, since duties derive from roles, promises, and relations and not from concern alone. From duty it does not follow that one is able, as the defunded official shows. And ability, finally, entails neither of the other three. An account of moral life that treats any of these as sufficient for any other will misdescribe numerous cases. The particular error that information networks encourage is the inference from knowledge, and the care it provokes, to duty and ability. The inference is plausible in the village, where the person who knows of a neighbor's trouble is often the person best placed to help. It fails at scale, where the person who knows is a member of an audience of millions.

\section{The Problem of Nonactionable Concern}

Responsibility without jurisdiction is a condition in which a person experiences pressure to respond to a problem without having adequate authority, information, or effective means of intervention. The phrase is meant to name a structure and not to excuse. It identifies a gap between the demand experienced and the capacity possessed, and it asks what follows from the existence of that gap.

Formally, let $\mathrm{Pressure}(a,p)$ hold when $a$ experiences a demand to respond to $p$, and let $\mathrm{Jur}(a,p)$ hold when $a$ possesses authority, information, and means adequate to intervene. Then agent $a$ is in a condition of responsibility without jurisdiction with respect to $p$ if and only if $\mathrm{Pressure}(a,p)\wedge\neg\mathrm{Jur}(a,p)$. The definition makes no claim about what the agent ought to do. It makes a claim about the structure of the situation.

It is essential to avoid suggesting that limited capacity automatically eliminates moral obligation. That a person cannot single-handedly end a war does not entail that she has no obligations concerning it. She may have obligations to inform herself responsibly, to support institutions that are capable, to refrain from contributing to the harm, to vote, to donate, and to refuse complicity. What limited capacity eliminates is the particular obligation to resolve, and it leaves in place obligations of a more modest and more appropriate kind. A principle that required every person to address every problem would be unsatisfiable, and an unsatisfiable principle is not a demanding principle but an empty one: it cannot guide action, because no action satisfies it. The ethical question is not whether to be concerned but what the appropriate form of the concern is for a person in a given position, and the thesis of the chapter is that the form is determined partly by jurisdiction.

\section{The Economy of Symbolic Participation}

Commentary, endorsement, condemnation, reposting, and numerical expressions of support can function as substitutes for material intervention. The mechanism is easily stated. A symbolic act is inexpensive, immediate, visible to others, and provides the actor with a sense of having responded. A material act is expensive, slow, frequently invisible, and uncertain in effect. A person facing pressure to respond who can relieve the pressure cheaply will often do so.

A minimal model clarifies the conditions under which substitution is complete. Let an agent require a felt contribution of at least $f_0>0$, where felt contribution from material effort $m\ge0$ and symbolic effort $s\ge0$ is $f=m+\theta s$ for some $\theta\in(0,1]$ representing the extent to which symbolic action is experienced as equivalent to material action. Let material effort cost $c_m$ per unit and symbolic effort cost $c_s$ per unit. The agent minimizes cost $c_mm+c_ss$ subject to $m+\theta s\ge f_0$.

\begin{proposition}[Complete Substitution]\label{prop:substitute}
The cost-minimizing agent chooses $m=0$ and $s=f_0/\theta$ whenever $c_s/\theta<c_m$, chooses $m=f_0$ and $s=0$ whenever $c_s/\theta>c_m$, and is indifferent when equality holds. In particular, if symbolic action is cheap enough that $c_s<\theta c_m$, no material effort is undertaken.
\end{proposition}

\begin{proof}
The problem is a linear program with a single constraint, and the binding constraint is $m+\theta s=f_0$. Substituting $m=f_0-\theta s$ with $0\le s\le f_0/\theta$ gives cost $c_mf_0+(c_s-\theta c_m)s$, which is linear in $s$ and is minimized at $s=f_0/\theta$ if the coefficient is negative, at $s=0$ if positive.
\end{proof}

The model is intentionally crude, and it is offered as a statement of conditions and not as an explanation of any actual behavior. It shows that complete substitution requires only that symbolic action be sufficiently cheap relative to the extent it is felt as equivalent, and it indicates why the proliferation of costless symbolic channels can reduce material action without any change in the intensity of concern. Whether any significant fraction of the population behaves in accordance with the model is an empirical question. The relevant literature on moral licensing offers mixed results, and the essay does not rely on the hypothesis that substitution is common.

The point should not be pressed beyond its warrant, because it is equally clear that communication and coordination are genuinely useful. Publicity can bring a hidden wrong to the attention of those who can remedy it. Collective expression can signal the existence of a constituency and shift the incentives of those in authority. Information shared among those with a stake can enable mutual aid. The criticism is directed at the performative case, in which the act has no plausible route to a consequence and its function is to signal the actor's stance. The line between the two is not always clear, and an act can be both, but the distinction is real: genuinely useful communication is addressed to someone who can act and carries information that bears on the action, and purely performative participation is addressed to an audience of equals who cannot.

\section{The Moral Costs of Perpetual Exposure}

The chapter would be incomplete without a discussion of the costs that perpetual exposure to distant problems may impose, and it is important to treat these as hypotheses requiring empirical evidence and not as universal consequences. Four are worth recording. Attentional exhaustion is the depletion of the capacity for sustained concern by the constant arrival of occasions for it, so that the concern available for the problems one can address is diminished by the concern spent on those one cannot. Misplaced confidence is the tendency to believe that one understands a distant matter on the strength of a stream of fragmentary reports, and to act on the belief. Shallow judgment is the substitution of rapid, affect-laden verdicts for the slower assessment that a complex matter requires. And the displacement of local obligations is the diminished attention to the commitments that arise in one's own household, workplace, and community, which are typically less dramatic and less rewarded with recognition than distant causes.

Each of these hypotheses has plausible mechanisms and some supporting testimony. Each is also subject to the objection that people differ greatly in how they respond, that some find in the stream of distant events a source of energy and engagement, and that the effects may depend heavily on how exposure is structured. The essay's argument does not require that the costs be universal. It requires only that they be real for enough people, in enough settings, that an information environment should be evaluated in part by whether it produces them, and that the burden of showing otherwise should not fall solely on those who experience the costs.

\chapter{The Disappearance of the Local Problem}

\section{What Makes a Problem Local?}

Locality should be defined relationally and not exclusively geographically. A problem is local to an agent when its relevant conditions are sufficiently accessible, its causal structure can be investigated, and the agent possesses some capacity to intervene and evaluate the consequences. Each clause carries weight. Accessibility concerns whether the agent can observe what is happening. Investigability concerns whether the agent can form and test hypotheses about why. The capacity to intervene concerns whether there is anything the agent can do, and the capacity to evaluate concerns whether the agent can tell afterward whether it worked.

On this definition geographic proximity is neither necessary nor sufficient. A geographically distant software defect may be more local to a developer than a dispute occurring on the same street. The developer can read the code, reproduce the failure, change the source, and run the tests. The resident of the street may know that the dispute is occurring, but may have no standing in it, no access to the parties' reasons, and no way of changing the outcome. What the definition captures is the degree to which an agent stands in a position of practical efficacy with respect to a problem, which is the position that permits learning from action, and geographic proximity is merely one of several things that can produce it. The definition also makes locality a matter of degree, and a problem can be more or less local to an agent along each of its dimensions.

\section{The Leaking-Pipe Example}

The repair of a leaking pipe serves as a paradigmatic example of bounded intervention. The problem has a discoverable location, since water appears where it should not. It has an identifiable failure mechanism, such as a cracked joint, a corroded section, or a failed seal. There is an available repertoire of repairs, including tightening, replacement, and sealing. And there are observable criteria for success, which are that the water no longer appears and that the fitting holds under pressure. A person with modest skill can move through these stages and arrive at a verified result in an afternoon.

The example should not be idealized, and it is more instructive if its difficulties are noticed. The visible leak may be the symptom of a failure elsewhere, as when a drip at a fitting results from excessive pressure upstream. Hidden water damage may have been proceeding behind a wall for months. An incorrect diagnosis may lead to the replacement of a sound part while the cause remains in place. A repair that appears successful under test may fail under conditions that the test did not reproduce. And a repair at one point may transfer stress to another, so that the downstream plumbing fails. Even apparently simple local problems can contain uncertain consequences, and the point of the example is not that local problems are easy. It is that they are tractable in a specific sense: the uncertainty can be reduced by looking, testing, and trying again, because the problem sits where the agent can reach it.

\section{From Observation to Correction}

The sequence by which a problem is corrected can be written as
\[
\text{Observation}\rightarrow\text{Diagnosis}\rightarrow\text{Intervention}\rightarrow\text{Verification}.
\]
Each stage can fail, and reliable correction requires more than merely detecting an undesirable condition. Observation fails when the condition is unnoticed or noticed in a distorted form. Diagnosis fails when the observed condition is attributed to the wrong cause. Intervention fails when the action taken does not accomplish what was intended, or accomplishes it with side effects. Verification fails when the evidence for success is inadequate, so that a failed correction is taken for a successful one or a successful correction is not recognized.

The structure supports a simple quantitative observation. Let $q_1$ be the probability that the observation stage succeeds, $q_2$ the conditional probability that diagnosis succeeds given a successful observation, $q_3$ the conditional probability that intervention succeeds given a successful diagnosis, and $q_4$ the conditional probability that verification succeeds in correctly confirming a successful intervention. By the chain rule of probability, the probability of a verified correction is the product $q_1q_2q_3q_4$. This number is no greater than the smallest factor and is generally much smaller than any of them when several are moderate. The practical implication is that a system which improves the first stage only, by making problems more visible, can raise the product by at most a factor of $1/q_1$, and cannot raise it at all if the bottleneck lies elsewhere. Visibility without diagnostic competence, the means of intervention, and the capacity for verification does not produce correction. It produces awareness.

\section{The Value of Bounded Problems}

Sufficiently bounded problems make it possible to learn from intervention. They permit causal hypotheses to be revised against observable outcomes, and this property is the source of the practical knowledge that distinguishes the experienced practitioner from the novice. The point can be formalized by a Bayesian illustration. Suppose that an agent repeatedly applies an intervention of unknown efficacy $\vartheta\in[0,1]$, with outcomes that are independent successes with probability $\vartheta$ and that can be observed. Starting from a uniform prior, after $n$ observed trials with $s$ successes the posterior is the Beta distribution with parameters $1+s$ and $1+n-s$, with mean $(1+s)/(2+n)$ and variance bounded above by $1/(4(n+3))$. The uncertainty about $\vartheta$ shrinks at rate $1/n$ in variance. If, by contrast, the outcomes of the intervention are never observed, the posterior equals the prior, and the variance remains $1/12$ however many times the intervention is performed. The agent in the second situation acts and learns nothing.

The difference is the difference between a feedback-rich and a feedback-poor environment, and the argument for locality is in large part an argument about feedback. Where the effects of an action are visible to the actor, practice accumulates into knowledge, and the actor's judgments improve. Where they are not, practice can accumulate into confidence without accumulating into knowledge, which is a more dangerous condition than ignorance. A civilization that organizes the attention of its members around matters from which no feedback returns is deliberately producing confident actors who cannot learn.

\section{The Danger of Abstraction Without Feedback}

The comparison is with engagement in controversies whose relevant evidence, institutional authority, and causal consequences remain inaccessible to the observer. A person who forms and expresses an opinion about a distant legislative dispute, an international conflict, or a technical controversy in a field outside her competence is engaged in an activity of a different kind from the person who repairs a pipe. She may inform herself well or badly. She may reach a conclusion that is right or wrong. But she cannot observe what her having reached it changes, and the activity therefore lacks the corrective structure of practical diagnosis.

The argument concerns the limitations of ungrounded intervention, and not the intellectual illegitimacy of studying distant events. Study is valuable. The scholar who analyses a distant conflict, the journalist who reports it, the analyst who models it are producing knowledge that others can use, and their activity has its own feedback in the form of criticism, replication, and correction by peers. The danger lies in the intermediate position of the person who has been given the impression of participation without being given the means of it. For that person the information environment supplies a stream of occasions to feel and to express and no occasion to test.

\part{Toward a Mind Your Own Business Protocol}

\chapter{Locality as a Principle of Agency}

\section{Defining the Protocol}

The Mind Your Own Business Protocol, abbreviated MYOB, is introduced here as a proposed allocation rule for attention, intervention, and responsibility. The name is deliberately plain, and it is chosen to be provocative in a particular way. Read as a command to ignore the troubles of others, the phrase would be a moral failure. The essay reads it otherwise, as a theory of bounded responsibility, legitimate intervention, and the allocation of attention. Its sense is that each agent has a business, a domain in which she has standing and capacity, and that the first task of practical ethics is to attend well to that domain, to be accountable within it, and to refer what lies beyond it to those whose business it is. Minding one's business in this sense is not the abdication of concern for others. It is the discipline of directing concern through the channels in which it can be effective.

The central presumption is that agents should prioritize problems for which they possess meaningful observational access, causal competence, legitimate authority, and reliable feedback. The presumption is defeasible, as Section 8.6 specifies, and it is a presumption about priority and not about exclusion. It does not say that agents should attend only to such problems. It says that when attention and effort must be allocated, as they always must, the problems that satisfy these four conditions have a prima facie claim to precedence over those that do not.

\section{The Four Dimensions of Actionable Intervention}

Define $O_i(p)$ as the observational access of agent $i$ to problem $p$, $C_i(p)$ as intervention capacity, $F_i(p)$ as feedback reliability, and $L_i(p)$ as legitimacy or authorization. A provisional model of intervention readiness is
\[
R_i(p)=O_i(p)\,C_i(p)\,F_i(p)\,L_i(p),
\]
with normalized terms in $[0,1]$. This is a heuristic measure of practical intervention readiness, not a complete theory of moral responsibility. It does not say what an agent owes, and it does not say what an agent may do. It says how well placed the agent is to do something effective and accountable.

The multiplicative form models bottlenecks, but the precise functional form must be justified and compared with alternatives. Some properties can be proved directly.

\begin{proposition}[Properties of the Readiness Index]\label{prop:readiness}
Let $x_1,\dots,x_4\in[0,1]$ be the four components and $R=\prod_jx_j$. Then
(a) $R=0$ if and only if some component is zero;
(b) $(\min_jx_j)^4\le R\le\min_jx_j$;
(c) the elasticity $\partial\ln R/\partial\ln x_j=1$ for every $j$ with $x_j>0$, so a given proportional improvement in any component has the same proportional effect on $R$;
(d) no additive index $S=\sum_jw_jx_j$ with positive weights summing to one can satisfy (a), since $S>0$ whenever any component is positive.
\end{proposition}

\begin{proof}
Part (a) holds because a product of real numbers vanishes if and only if a factor does. For (b), each $x_j\ge\min_kx_k$ gives $R\ge(\min_kx_k)^4$, and since the remaining factors are at most one, $R\le x_j$ for each $j$, hence $R\le\min_jx_j$. For (c), $\ln R=\sum_j\ln x_j$, whose partial derivative with respect to $\ln x_j$ is one. For (d), if $x_1>0$ and the other components are zero then $S\ge w_1x_1>0$.
\end{proof}

Part (d) is the reason for the multiplicative form: a missing authorization should not be compensable by abundant capacity, and the additive index would let it. Part (b) shows that the product is at least as demanding as the minimum, which may be too demanding: two moderate deficits depress $R$ more than one severe deficit would under a pure minimum rule. The geometric mean $R^{1/4}$ is a natural alternative that retains the veto property (a) while being less severe, and a still different alternative is to treat legitimacy not as a factor but as a constraint, admitting only agents with $L_i(p)=1$ into the comparison. The choice among these is a normative decision and not a mathematical one, and the essay does not decide it. What is argued is that any adequate index must have the veto property for the components that represent preconditions.

\section{Competence Is Not Entitlement}

An agent may be capable of changing a system without possessing legitimate authority to do so. A person with the technical skill to alter the records of an institution is not thereby authorized to do so, and a person who can see how a neighbor ought to arrange his household has no standing to arrange it. The protocol must therefore distinguish practical ability from permission and obligation. This is why legitimacy appears as a separate factor in the readiness index, and why its zero value vetoes the index regardless of the value of the others.

The distinction guards against a familiar pathology of technical competence, which is the belief that the ability to fix a thing creates the right to fix it. The history of technocratic intervention contains many instances of competent persons imposing solutions that were effective in the sense of the problem as they defined it and objectionable in the sense of the persons they affected. Legitimacy is not a formality that the competent may waive. It is the recognition that the problem is also someone else's, and that the people to whom it belongs have claims on how it is addressed. Section 8.5 develops the corresponding positive virtue.

\section{Locality and Subsidiarity}

The proposal can be compared with the principle of subsidiarity in political and institutional theory. In the social teaching of the Catholic Church, notably in the encyclical \emph{Quadragesimo Anno} \cite{pius}, the principle holds that a higher body should not take over functions that a lower body can perform adequately, and that the higher body's proper role is to support the lower. The principle was later incorporated in the constitutional law of the European Union, where Article 5 of the Treaty on European Union \cite{tue} states it as a criterion for the exercise of non-exclusive competences. The key question is whether decisions should ordinarily be made at the lowest level capable of handling them competently, with escalation when local resources or authority are inadequate.

Subsidiarity and the protocol share this structure, and differ in their subject. Subsidiarity concerns the allocation of authority among levels of an organization. The protocol concerns the allocation of an individual's attention and effort among problems. The two are connected, since the assignment of authority to a level is in part an assignment of the attention that level will give. Both rest on the observation that the lower level typically has better information about the matter, is more directly affected by the outcome, and can be held to account more readily. Both also admit the qualification that the lower level can fail, and that the higher level may then have a duty to act. The comparison suggests that the protocol is not an eccentric proposal but an application to individual conduct of a principle that has a long history in institutional design.

\section{Nonintervention as a Positive Capacity}

The argument is not only that agents should do what they can, but that knowing when not to intervene is itself a form of practical competence. Restraint may preserve the agency of others, prevent interference, and reduce the consequences of insufficient information. A physician who declines to prescribe when the evidence for benefit is thin, a manager who lets a team resolve its own disagreement, and a bystander who recognizes that a scene she does not understand is not hers to resolve are exercising competence, not displaying its absence.

The case for restraint is epistemic as well as moral. An intervention based on partial information has a risk of harm that increases with the ignorance of the intervener, and the risk is often not visible to the intervener. A decision not to act because one lacks the information to act well is a decision responsive to evidence. This is the opposite of indifference. It is the application of the same standard of care to the question whether to act as to the question how. The positive virtue is the ability to distinguish a situation that calls for action from one that merely invites it, and the ability is acquired, like other competences, by experience with feedback.

\section{A Defeasible Presumption, Not an Absolute Prohibition}

The protocol creates a default preference for local intervention, subject to exceptions involving external harms, emergencies, institutional duties, and failures of local accountability. The structure is that of a default rule in the sense of nonmonotonic logic. Let $\mathrm{Pref}(i,p)$ denote that agent $i$ gives $p$ priority, and let $\theta$ be a readiness threshold. The default is
\[
R_i(p)\ge\theta\ \wedge\ \neg\mathrm{Ex}(i,p)\ \rightsquigarrow\ \mathrm{Pref}(i,p),
\]
where the exception predicate is
\[
\mathrm{Ex}(i,p)\ \equiv\ \mathrm{ExtHarm}(p)\vee\mathrm{Emerg}(p)\vee\mathrm{InstDuty}(i,p)\vee\mathrm{LocalFail}(p),
\]
that is, an external harm falling beyond the local horizon, an emergency, an institutional duty arising from a role, or a failure of local accountability. When the exception holds, the default is defeated and a different consideration governs. The point of stating the rule in default form is to make clear that it is a presumption which shifts the burden of justification, and that the burden is met by establishing an exception. An agent who departs from the default must be prepared to say which exception applies, and an agent who adheres to it must be prepared to say why none does. Chapter 10 treats these exceptions in detail, since they are the substance of the objection that locality amounts to moral isolation.

\chapter{The Mathematics of Attentional Jurisdiction}

\section{A Finite Attention Budget}

Let agent $i$ have a finite attentional budget $B_i>0$, distributed among problems $p\in\cP$. Let $a_{ip}\ge0$ represent attention allocated by $i$ to $p$, subject to
\[
\sum_{p\in\cP}a_{ip}\le B_i .
\]
Attention is here a scalar resource, a simplification that suppresses the structure of attention in time, its dependence on context, and the cost of switching. These omissions are acknowledged and the scalar model is retained because it is the simplest in which the economics of allocation can be displayed.

\section{Expected Corrective Benefit}

Define $G_{ip}(a_{ip})$ as the expected benefit of allocating attention $a_{ip}$ to problem $p$, conditional on information quality, intervention capacity, legitimacy, and possible adverse consequences. The function incorporates the factors of the readiness index, in the sense that a problem for which $i$ has no capacity or no legitimacy yields small or zero benefit however much attention is given. The agent's allocation problem becomes
\[
\max_{\{a_{ip}\}}\ \sum_{p\in\cP}G_{ip}(a_{ip})\quad\text{subject to}\quad \sum_pa_{ip}\le B_i,\ a_{ip}\ge0 .
\]
For the remainder of the chapter the agent's index $i$ is suppressed.

\section{Opportunity Costs}

Allocating attention to one problem reduces the capacity to address another. This truism is the source of the main point. The cost of universalized attention is not simply fatigue. It includes the loss of corrective opportunities that were already within the observer's competence. A person whose attention is spread across a thousand distant problems has withdrawn attention from the few that she could have corrected, and has done so without any compensating increase in the correction of the others, since the marginal benefit of her attention to them is negligible. The loss is invisible, because the opportunities forgone do not appear on any ledger.

\section{An Attention-Allocation Proposition}

Under concavity and differentiability assumptions, the familiar condition for an optimal interior allocation is that marginal expected benefit be equalized across eligible problems.

\begin{proposition}[Equimarginal Allocation]\label{prop:equimarg}
Suppose each $G_p:[0,\infty)\to\RR$ is concave and differentiable. A feasible allocation $a^\ast$ with $\sum_pa^\ast_p=B$ is optimal if and only if there exists $\lambda\ge0$ such that $G_p'(a_p^\ast)=\lambda$ whenever $a_p^\ast>0$ and $G_p'(a_p^\ast)\le\lambda$ whenever $a_p^\ast=0$.
\end{proposition}

\begin{proof}
For sufficiency, take any feasible $a$. Concavity gives $G_p(a_p)\le G_p(a_p^\ast)+G_p'(a_p^\ast)(a_p-a_p^\ast)$. Summing, $\sum_pG_p(a_p)\le\sum_pG_p(a^\ast_p)+\sum_pG_p'(a^\ast_p)(a_p-a^\ast_p)$. For $p$ with $a_p^\ast>0$ the derivative equals $\lambda$, and for $p$ with $a^\ast_p=0$ the term is $G_p'(a_p^\ast)a_p\le\lambda a_p$. Therefore the last sum is at most $\lambda\sum_p(a_p-a_p^\ast)\le0$ because $\sum_pa_p\le B=\sum_pa^\ast_p$. For necessity, suppose $a^\ast$ is optimal. If two problems $p,q$ with $a_p^\ast>0$ had $G_p'(a_p^\ast)<G_q'(a_q^\ast)$, moving a small amount of attention from $p$ to $q$ would raise the sum, a contradiction, which forces equality of derivatives among the positively funded problems. If $a_q^\ast=0$ with $G_q'(0)$ exceeding the common value $\lambda$ on funded problems, moving attention to $q$ would likewise increase the total. Finally, $\lambda\ge0$ holds when the budget constraint binds because otherwise reducing allocation to any funded problem would not lower the sum.
\end{proof}

The result is conditional on an adequate representation of harms, obligations, and uncertainty. The functions $G_p$ summarize the entire causal structure of the interaction between attention and benefit, and the equimarginal rule is no better than the specification of $G_p$. It says nothing about which problems are worth solving, only how to divide a fixed budget once their benefit functions are given. The condition is nevertheless informative, since it shows that the correct criterion for allocating attention is the marginal benefit of an additional unit and not the size, visibility, or urgency of the problem. A large and highly visible problem whose benefit function is flat at the margin deserves no more attention than a small one.

The comparison with allocation according to visibility can be made exact in the quadratic case. Let $G_p(a)=b_pa-\tfrac12c_pa^2$ with $b_p>0$ and $c_p>0$, and let the budget be large enough that the optimal allocation is interior in the sense of assigning positive attention to every problem.

\begin{proposition}[Local Corrective Allocation]\label{prop:local}
In the quadratic case, the optimal allocation is $a_p^\ast=(b_p-\lambda)/c_p$ with $\lambda=\bigl(\sum_qb_q/c_q-B\bigr)\big/\sum_q(1/c_q)$, provided all $a^\ast_p>0$. For any feasible allocation $a=a^\ast+\delta$ with $\sum_p\delta_p=0$ the loss in total benefit relative to the optimum is exactly
\[
\sum_pG_p(a_p^\ast)-\sum_pG_p(a_p)=\tfrac12\sum_pc_p\,\delta_p^2 .
\]
In particular, allocation in proportion to visibility, $a_p^v=B\,v_p/\sum_qv_q$, achieves the optimum if and only if $a^v=a^\ast$, and otherwise incurs the loss $\tfrac12\sum_pc_p(a_p^v-a_p^\ast)^2>0$. The visibility allocation can equal the optimum in special cases, for example if $c_p=c$ for all $p$ and $b_p-\lambda$ is proportional to $v_p$; and an allocation based on estimated benefits may be worse than the visibility allocation if the estimates deviate from the truth by more, in the $c$-weighted norm, than visibility does.
\end{proposition}

\begin{proof}
The first-order conditions $b_p-c_pa_p=\lambda$ give $a_p=(b_p-\lambda)/c_p$, and imposing $\sum_pa_p=B$ determines $\lambda$ as stated. For the loss, expand $G_p(a^\ast_p+\delta_p)=G_p(a^\ast_p)+G_p'(a_p^\ast)\delta_p-\tfrac12c_p\delta_p^2$, which is exact for a quadratic. Since $G_p'(a_p^\ast)=\lambda$ for every $p$, the linear terms sum to $\lambda\sum_p\delta_p=0$, leaving the quadratic loss. The statements about the visibility allocation follow by substituting $\delta=a^v-a^\ast$, and the final claims are direct consequences of the formula: the loss of any allocation is its $c$-weighted squared distance from the optimum, and an estimate-based allocation beats the visibility allocation exactly when it is closer to the optimum in that norm.
\end{proof}

The proposition is a precise form of the claim that allocation according to expected corrective benefit can outperform allocation according to visibility or message frequency. The benchmark is defined: the visibility rule allocates in proportion to a nonnegative visibility index. The counterexample is also supplied, in the sense that the claim is false when visibility happens to be aligned with marginal benefit, and the proposition is silent about which situation is more common. The claim is also conditional on the quality of the benefit estimates. The quadratic loss formula implies that an allocation rule is only as good as its closeness to the optimum, and an agent who has no reliable estimate of benefit may do better to follow visibility, which at least reflects the aggregate judgment of many observers. The argument for the protocol is therefore not that estimated benefit always dominates visibility, but that where reliable estimates of local benefit are available, as they are for problems that satisfy the four conditions, they should be used in preference.

\section{The Problem of Unmeasured Value}

The proposed optimization itself risks reproducing Lanier's problem. If the protocol evaluates every activity using a numerical score, it may create a new ticker devoted to maximizing actionable responsibility. A person who has internalized the equimarginal rule might feel obliged to compute the marginal benefit of each moment, to track her own allocation against the optimum, and to reproach herself for any deviation. A score of responsibility is a metric like any other, and the analyses of Chapters 3 and 4 apply to it. It would be subject to Goodhart's law, would acquire a weight $\alpha_t$ that increases with exposure, and would be open to gaming.

This self-referential difficulty must be made explicit, because the philosophical proposal and the mathematical apparatus pull in different directions. The apparatus is useful for analysis, for the design of institutions, and for the evaluation of allocation rules. It is not a procedure that an individual should run continuously in the course of living. The model describes the structure of the problem from the standpoint of a designer, and is not a guide to the moment-to-moment conduct of an agent. A person who adopts the protocol should not be maximizing anything. She should be following a set of boundaries.

\section{Escaping Compulsory Optimization}

The alternatives to continuous maximization are coarse thresholds, qualitative judgments, bounded decision procedures, and satisficing. Satisficing, introduced by Herbert Simon \cite{simon56} as a model of choice by agents with bounded computational capacity, replaces the search for the best option by the search for one that meets an aspiration level and the termination of search when it is found. Applied to attention, it replaces the question of how to allocate each unit optimally by the question of whether the problems within one's jurisdiction have received enough attention to count as handled. The aspiration levels can be set by role, by norm, and by explicit agreement, and they provide the condition of closure that continuous maximization lacks.

The protocol should help agents determine when sufficient attention has been given, and not merely encourage them to optimize a new measure indefinitely. The distinction between a rule that tells one when to stop and a rule that tells one how to do better is fundamental. The former defines a task that can be finished, and the latter defines a task that cannot. Chapter 18 builds on this distinction in stating the right to finish. For present purposes, the mathematical result of this chapter should be read as a result about the design of an institution that wishes to assist agents in avoiding misallocation, and its use as a personal calculus should be resisted.

\chapter{Mind Your Own Business Is Not Moral Isolation}

\section{The Externality Objection}

The most immediate objection to the protocol is that it appears to license moral isolation. The cases are familiar. Pollution imposes costs on those downstream who cannot see its source. Financial fraud harms victims who do not know the perpetrator and cannot observe the instruments. Public health depends on the collective conduct of persons who individually observe only their own health. Environmental degradation proceeds by accumulation beyond any individual's horizon. Institutional abuse flourishes where those inside are silent and those outside are uninformed. In each of these, harms extend beyond the immediate observational horizon of the persons who cause or could prevent them, and a principle that directs agents to attend to what they can observe appears to direct them away from precisely the matters that most need attention.

The objection has force, and the protocol does not evade it by denying the facts. It responds by observing that the structure of the objection is an argument for institutions. If a harm is invisible to individuals, then individuals cannot be expected to address it by attending to it, and the appropriate response is to create an institution whose function is to make it visible and to assign the responsibility for correcting it. The failure the objection identifies is a failure of the local principle to be supplemented by the means of extending locality, not a failure of the principle.

\section{The Limits of Immediate Perception}

A person may be responsible for consequences that she cannot directly observe. The absence of immediate feedback cannot excuse negligence where foreseeable harms can be investigated through appropriate methods. A manufacturer is responsible for the safety of a product used by customers she will never meet, and the responsibility is discharged by testing, by monitoring reports of failure, and by responding to them. A landowner who stores toxic material is responsible for its escape into the water table even though the escape will become evident only at a distance and after delay. The foreseeability of the harm and the availability of methods of investigation define the scope of duty, regardless of the perceptual horizon of the agent.

This qualification shows that the protocol's emphasis on observation should be understood as an emphasis on observability, which is a property of the problem and the available methods and not of unaided perception. A problem can be observable to an agent through instruments, records, reports, and audits though it lies far beyond her senses. The question the protocol poses is whether the agent can, through some accessible means, determine what is happening and what her actions contribute to it. Where she can and fails to, she is negligent. Where she cannot, the matter is properly assigned elsewhere.

\section{Extending Locality Through Institutions}

Scientific instruments, legal procedures, professional expertise, and independent audits can extend the effective observational and corrective reach of individuals. A thermometer extends the sense of heat to a precision and in a range that the skin cannot supply. A court extends the capacity to adjudicate by providing procedures through which a dispute between strangers can be resolved with authority. A professional license extends the trust that a client can place in a practitioner whose competence she cannot assess for herself. An audit extends the observation of an organization's conduct to those who cannot enter it. Each of these is a means of making a nonlocal matter local, in the relational sense of Chapter 7.

Locality is therefore not a fixed geographic boundary but an institutional accomplishment. The history of public health illustrates the point. Before the identification of the mechanisms of waterborne disease, the contamination of a water supply was not observable by any individual in any useful way, and a citizen could not be held responsible for what she could not detect. The development of epidemiological methods, sanitary engineering, and regulatory agencies converted an unobservable matter into an observable one and assigned responsibility to those equipped to act. The principle of locality was not violated in the process. It was satisfied by extending the range over which its conditions held.

\section{Escalation and Delegation}

The protocol requires procedures for referring problems to competent agents and not broadcasting them indiscriminately to everyone. A person who becomes aware of a problem beyond her own competence has several options besides acting herself or ignoring it. She may report it to a body responsible for it. She may supply information to those who can verify it. She may support an organization that works on it. She may place it in the hands of a more competent acquaintance. Each of these is a form of delegation, and each is an exercise of responsibility. The distinction between delegation and broadcast is that delegation identifies a recipient who can act and ensures the problem arrives with the information needed to act, while broadcast distributes the problem to a large audience in the hope that some member is capable. Broadcast has its uses, as when the identity of the competent party is unknown, but it should be a fallback and not a default, and when used it carries a cost that Chapter 5 quantified: the diffusion of the problem across many who cannot help dilutes its reach to those who can.

\section{The Ethical Distinction Between Concern and Interference}

Legitimate concern can motivate investigation, support, delegation, or institutional reform without implying that every concerned observer should intervene directly. The distinction between concern and interference corresponds to a familiar distinction between caring about a matter and taking charge of it. A person may care deeply about the conduct of a distant government and express that care by informing herself, voting, funding relevant organizations, and supporting accountability mechanisms, without any suggestion that she should take direct action against officials. A person may care about the well-being of a neighbor and express it by being available and attentive, without taking over the neighbor's affairs. The protocol concerns the form of caring appropriate to a position, and it endorses concern as the condition from which responsible action starts. What it opposes is the conversion of concern into a demand for direct intervention by those unequipped for it, and the corresponding conversion of every observer into a claimed participant.

\part{Institutions, Incentives, and the Reproduction of Problems}

\chapter{When Institutions Are Rewarded for Activity Rather Than Correction}

\section{The Institutional Ticker}

The analysis can be extended from individual interfaces to organizational dashboards. An institution, like an individual, is surrounded by counters. It may count cases processed, inspections completed, interventions delivered, reports published, or funds distributed, and these counts may be displayed to its leadership, its funders, and the public. Such counts can be administratively useful without establishing that the underlying problem has diminished. A department that processes a larger number of cases may be processing them because the number of cases has grown, because it has become more efficient, or because it has lowered the threshold at which a matter is counted as a case. The count is silent among these alternatives, and a dashboard that displays only the count invites the inference most flattering to the institution.

The institutional ticker has the same structure as the individual one. It is a continuously updated representation whose motion is generated by the activity of the organization and whose relation to the state of the world is indirect. The dashboard rises when work is done. Whether the work corrects anything is a further question that the dashboard does not answer, and in many institutions it is a question that no one is assigned to ask.

\section{Activity Metrics and Outcome Metrics}

Four categories of measure should be distinguished. Workload measures the demand placed on the institution, such as the number of requests received. Output measures what the institution produces, such as the number of cases closed or services delivered. Outcome measures the condition of the people or systems the institution is meant to serve, such as the incidence of the harm it addresses. And verified reduction of harm measures the portion of any change in outcome that can be attributed, with independent evidence, to the institution's action. The categories are ordered by the difficulty of measuring them and by their closeness to the purpose of the institution, and the ordering is inverse: the measures easiest to obtain are furthest from the purpose.

Because administrative systems naturally record workload and output, and because outcomes and attributable reductions require additional effort, independent data, and often years, institutions systematically report the former. This is not usually a matter of deception. It is a consequence of what is cheap to count. The result is that institutional self-description gravitates toward the measures most remote from correction, and the gravitation is reinforced when workload and output are the quantities on which funding and reputation depend.

\section{The Correction-Reproduction Problem}

Organizations may acquire incentives to preserve the problems that justify their funding, authority, or continued existence. The idea is old and has been expressed in various forms, from the observation that a bureaucracy tends to expand its own mandate to the more pointed claim that a charity that eliminated its cause would eliminate itself. The essay does not assume that every institution exhibits this pathology. Most do not exhibit it in a severe form, and many are staffed by people whose commitment to the mission exceeds their interest in the institution's survival. The task is to specify the incentive arrangements under which the pathology becomes plausible, so that those arrangements can be recognized and avoided.

A simple model displays them. An institution faces an initial stock of problem $S_0>0$ and chooses a correction effort $e\in[0,S_0]$ at cost $\tfrac12ce^2$, leaving a residual problem $S=S_0-e$. Society derives a benefit $b>0$ per unit of problem removed, so that social welfare is $be-\tfrac12ce^2$ and is maximized at $e^{\mathrm{soc}}=b/c$ when $b/c\le S_0$. The institution is funded according to some rule, and its objective is funding minus cost.

\begin{proposition}[Funding Rules and Correction Effort]\label{prop:funding}
Under need-based funding $F=\varphi S$ with $\varphi>0$, the institution's utility $\varphi(S_0-e)-\tfrac12ce^2$ is strictly decreasing in $e$ and the institution chooses $e=0$, correcting nothing. Under outcome-based funding $F=\varphi(S_0-S)=\varphi e$, the institution chooses $e=\min(\varphi/c,S_0)$, which equals the socially optimal effort when $\varphi=b$ and $b/c\le S_0$. Suppose now that the institution may also overstate its reduction by $m\ge0$ at cost $\tfrac12km^2$, where an independent verifier detects and disallows the overstatement with probability $\pi\in[0,1]$. Under outcome-based funding with undetected reports paid at the same rate, the utility is $\varphi\,(e+(1-\pi)m)-\tfrac12ce^2-\tfrac12km^2$, and the optimum is $e=\varphi/c$ and $m=(1-\pi)\varphi/k$. The overstatement is zero if and only if $\pi=1$ and decreases linearly in $\pi$.
\end{proposition}

\begin{proof}
For need-based funding the derivative of utility with respect to $e$ is $-\varphi-ce<0$ for all $e\ge0$. For outcome-based funding, the derivative is $\varphi-ce$, vanishing at $\varphi/c$, with the constraint $e\le S_0$ binding otherwise. In the third case the objective is strictly concave and separable in $e$ and $m$, with first-order conditions $\varphi-ce=0$ and $(1-\pi)\varphi-km=0$.
\end{proof}

The proposition identifies the conditions under which correction fails. Funding that increases with the residual stock of the problem pays for its persistence, and the institution that responds to incentives corrects nothing, even though correction is socially valuable. The remedy is to pay for verified reduction, and the cost of the remedy is the cost of verification, since the benefit of outcome-based funding is eroded by overstatement in proportion to the weakness of verification. The three ingredients of the pathology are therefore funding that rises with the measured size of the problem, the absence of independent verification of outcomes, and the absence of any mechanism that ends the institution's mandate when the problem is resolved. The first two are modeled; the third is considered below. Where all three hold, the incentive to preserve the problem is not a failing of character but the rational response to the structure.

\section{Independent Verification}

The role of independent verification is therefore central. Evaluators whose budgets, leadership, and rewards are not controlled by the organizations whose outcomes they assess are able to supply credible measures of outcomes. The parameter $\pi$ in the model is not exogenous: it is a function of the evaluator's independence. An evaluator funded by the organization evaluated has an interest in finding that the organization is succeeding, and the detection probability $\pi$ is correspondingly lowered. An evaluator whose reputation depends on the accuracy of her findings has an interest in detection. The structural conditions of independence are three. The evaluator's funding must not depend on the evaluated body's goodwill. The evaluator's personnel must not be drawn from, and not look forward to employment with, the evaluated body. And the evaluator's methods and data must be open to scrutiny by third parties, so that findings can be contested. The audit literature offers a cautionary qualification \cite{power}: audit can itself become a ritual in which the form of verification is performed without its substance, and an evaluation regime can be gamed in the way that the activity it evaluates is. Independent verification is therefore a necessary condition for the integrity of outcome measures and not a sufficient one, and it must be supported by transparency about methods.

\section{Mandate Review}

Periodic institutional reviews can assess whether the original problem persists, whether intervention remains effective, and whether responsibility should be transferred or the institution reorganized. The concept is that of a sunset provision applied to mandates: an institution's authority is justified by a problem, and the justification should be revisited when the problem changes. A review might ask four questions. Does the problem the institution was created to address still exist at a comparable scale? Is the institution's activity reducing it, as independently verified? Is the institution the appropriate body to continue to address it, or has the capability moved elsewhere? And what would the consequences be of winding down, merging, or redirecting the institution?

Care is required in the design of review because it can generate the pathology it seeks to prevent. If continued funding depends on the problem persisting, then a review that terminates the mandate when the problem is solved provides the institution with a strong incentive not to solve it. The incentive can be reduced by arrangements that make the successful completion of a mission rewarding for those who accomplished it, such as guaranteed redeployment of staff, recognition of the achievement, and transition support. The general principle is that the end of a mandate should not be punished for the people who bring it about. A review that threatens institutional death for success is a review that will be met with strategic failure.

\section{The Difference Between Ending Activity and Solving the Problem}

An institution can reduce its reported workload by improving conditions, suppressing complaints, refusing cases, or manipulating classification rules. A falling count of cases is compatible with all four, and the first alone is a success. A reduction in reported incidents may reflect safer conditions, but it may equally reflect a discouraged public who have stopped reporting, a narrowed definition of what counts as an incident, or a decision to close matters without resolution. The essay insists upon independent evidence of actual improvement. Where the only evidence of improvement is the institution's own count, the improvement is unproven, and an observer who accepts it is relying on the very measure whose reliability is in question.

This insistence is the institutional analogue of the right to finish stated in Chapter 18. An institution has finished its work when the problem is resolved and the resolution can be independently confirmed, and it has not finished merely because its dashboard has stopped moving or has begun to move downward.

\chapter{Relationship Costs and the Price of Harmful Participation}

\section{The Relationship-Cost Principle}

The idea that repeated participation in demonstrably harmful transactions can increase the cost of maintaining institutional relationships deserves examination, with attention to its dangers. The principle states that an agent who repeatedly engages in conduct that independent findings establish to be harmful to others should expect the terms on which others deal with her to worsen, through higher prices, stricter conditions, more intensive oversight, or in extreme cases exclusion. The principle is familiar in ordinary economic life. An insured driver with repeated accidents pays more, a borrower with defaults is charged higher interest, and a supplier with documented quality failures is audited more closely.

The illustrative case is the arms dealer, chosen because the trade is lawful in many of its forms and gravely harmful in others, so that a single occupation spans the whole range the principle must handle. Three kinds of transaction can be distinguished. Supply for defense, where the buyer is a state or group resisting an aggression and the transfer is justified and verifiable, carries a low relationship cost, since both the probability and the severity of harm attributable to the transaction are small. A gray-area transaction, in which the end use is uncertain, the buyer's conduct is disputed, or the goods are readily diverted, carries a moderate cost reflecting the higher risk. Supply to an aggressor, where independent findings establish that the buyer is using such goods to commit grave harms, carries a high cost, since both the probability of harmful consequence and its severity are high. Three situations in law and evidence should be kept apart as well: lawful trade, in which no wrongdoing is alleged; disputed conduct, in which allegations are made but not established; and established wrongdoing, in which independent findings have been made through a process that the affected party had the opportunity to contest. The relationship-cost principle properly applies only to the third and to the portions of the second that independent evidence has substantiated. Applying it to allegation alone would convert accusation into punishment, and applying it to lawful trade would convert disapproval into sanction. The distinctions are part of the principle and not external qualifications.

\section{Friction Rather Than Prohibition}

The principle is not a prohibition, and the difference matters for its design. Some activities should simply be prohibited, namely conduct that is unacceptable under any circumstances, where no benefit could justify the harm and where the harm is of a kind that pricing cannot compensate, such as irreversible mass injury. For other conduct, a higher cost can reduce harm without a total ban, by making the transaction proceed only when its benefits are worth the additional cost. The two instruments differ in what they presuppose. A prohibition presupposes that the harm is categorically unacceptable. A price presupposes that the harm is compensable in principle and that the question is one of degree, and where that presupposition is false a price is the wrong instrument, because it invites the wealthy to buy what ought not to be sold.

Make the pricing instrument precise. For a transaction between agents $i$ and $j$, let the relationship cost be
\[
\tau_{ij}=\lambda\,q_{ij}\,s_{ij},
\]
where $q_{ij}\in[0,1]$ is the probability of a harmful consequence attributable to the transaction, $s_{ij}\ge0$ is the estimated severity of that harm, and $\lambda>0$ is a parameter fixing how strongly risk is priced. The cost is an expected harm scaled by a pricing intensity, and it is added to the base price and the ordinary transaction cost. In the figure of dollar amounts that accompanies the original sketch of the idea, a deal with a base price of 100 million and ordinary transaction costs of 5 million carries a relationship cost of 30 million, for a total of 135 million. Those amounts are illustrative and carry no empirical weight. What the example shows is the logical structure: the higher cost makes the deal less attractive, and the deal can still occur, but only if its benefits are worth the additional cost.

\begin{proposition}[Friction and Prohibition]\label{prop:friction}
Suppose the parties to a transaction obtain a joint net benefit $b_{ij}\ge0$ before the relationship cost, and the transaction proceeds if and only if $b_{ij}\ge\tau_{ij}$. Then the set of transactions that proceed is monotonically shrinking in $\lambda$, in $q_{ij}$, and in $s_{ij}$; for every transaction with $q_{ij}s_{ij}>0$ there is a finite $\lambda^\ast=b_{ij}/(q_{ij}s_{ij})$ above which it does not proceed; and no finite $\lambda$ excludes a transaction with $q_{ij}s_{ij}=0$. Consequently, friction approximates prohibition for any bounded benefit but cannot substitute for it where the benefit is unbounded or the harm is not compensable.
\end{proposition}

\begin{proof}
The condition $b_{ij}\ge\lambda q_{ij}s_{ij}$ is weakened by decreasing any of the three factors, which gives monotonicity. Solving $b_{ij}=\lambda q_{ij}s_{ij}$ for $\lambda$ gives the threshold when $q_{ij}s_{ij}>0$, and when $q_{ij}s_{ij}=0$ the right side is zero for every $\lambda$, so the condition holds.
\end{proof}

\section{Pricing Harm Through the Network}

Relationship costs can apply directly, between the parties to a transaction, and in some cases indirectly through supply chains. The arms dealer sits at the center of a network of raw materials suppliers, shipping companies, equipment providers, financial institutions, and insurers, each of which contributes to the dealer's capacity to supply and each of which may face some exposure. The question is how far the cost should spread, and the danger in answering it is the pathology of guilt by association, in which costs spread indefinitely through the network and attach to parties whose contribution to any harm is remote.

The disciplined answer is that indirect exposure should decline with distance. Let $k$ denote the number of steps between a party and the harmful transaction in the network, and let a weight $w(k)=e^{-\alpha k}$ with $\alpha>0$ determine the fraction of the direct cost borne at distance $k$. A party at distance $0$ bears the full cost, and each further step multiplies the exposure by $e^{-\alpha}<1$. The geometric form is chosen for convenience and is one of several that would serve, and its justification is that it gives a bounded total.

\begin{proposition}[Bounded Exposure Under Distance Decay]\label{prop:decay}
Let each party at distance $k$ from a harmful transaction have at most $\beta^k$ counterparties at distance $k+1$ for some branching bound $\beta\ge1$, and let the cost borne at distance $k$ be $\tau e^{-\alpha k}$ per party. If $\beta e^{-\alpha}<1$, the total cost borne across the network is at most $\tau/(1-\beta e^{-\alpha})$, which is finite, and the share of the total borne beyond distance $K$ is at most $(\beta e^{-\alpha})^{K+1}$. If $\beta e^{-\alpha}\ge1$ the bound diverges, and costs can spread without limit.
\end{proposition}

\begin{proof}
The number of parties at distance $k$ is at most $\beta^k$ (taking the direct party at distance $0$ as one), so the total cost is at most $\tau\sum_{k\ge0}\beta^ke^{-\alpha k}=\tau\sum_{k\ge0}(\beta e^{-\alpha})^k$, a geometric series that converges to $\tau/(1-\beta e^{-\alpha})$ when the ratio is below one. The tail beyond $K$ is $\tau(\beta e^{-\alpha})^{K+1}/(1-\beta e^{-\alpha})$, so its share of the bound is $(\beta e^{-\alpha})^{K+1}$.
\end{proof}

The condition $\beta e^{-\alpha}<1$ is a design constraint with a clear reading: the decay rate must outpace the branching of the network. A dense network, in which each party deals with many others, requires a steeper decay to keep the total bounded, and a rule that ignores branching can produce unbounded exposure even when the decay at each step appears generous. This is the quantitative form of the safeguard against guilt by association. It is an argument for limiting the reach of indirect costs by an explicit rule and not by discretion.

\section{Internalizing External Costs}

The principle relates to several established mechanisms for making agents bear the external consequences of their conduct. Pigouvian taxation \cite{pigou} sets a charge on an activity equal to the marginal external damage it causes, so that the agent's private calculation reflects the social cost. The result can be stated precisely.

\begin{proposition}[Pigouvian Internalization]\label{prop:pigou}
Let an agent choose an activity level $q\ge0$ with private benefit $B(q)$, concave and differentiable, and let the activity impose external damage $D(q)$, convex and differentiable. Suppose the social optimum $q^\ast$ is interior and satisfies $B'(q^\ast)=D'(q^\ast)$. If the agent is charged $\tau=D'(q^\ast)$ per unit, the agent's optimal choice is $q^\ast$.
\end{proposition}

\begin{proof}
The agent maximizes $B(q)-\tau q$, a concave function, whose first-order condition is $B'(q)=\tau=D'(q^\ast)=B'(q^\ast)$. Strict concavity of $B$ makes $B'$ strictly decreasing and so $q=q^\ast$ is the unique solution.
\end{proof}

Liability regimes, insurance pricing, and risk-adjusted contracting achieve similar internalization by different routes. Liability makes the agent pay for damage after the fact, insurance prices the expected damage in advance and passes it to the agent as a premium, and risk-adjusted contracting writes the terms of a relationship so that they reflect the counterparty's history. The relationship-cost principle belongs to this family. Its distinctive feature is that the cost is imposed through the terms of continuing relationships, and not by a central authority, and therefore it relies on information about conduct being available to those who set the terms.

\section{Dynamic Costs}

Let $C_{ij}(t)$ represent the cost of maintaining a relationship between agents $i$ and $j$, adjusted according to independently verified conduct and relevant risk. Investigate how repeated harmful conduct could alter contractual terms or oversight requirements. A minimal dynamic is
\[
C_{ij}(t+1)=(1-r)\,C_{ij}(t)+\gamma\,h_{ij}(t),\qquad r\in(0,1],\ \gamma>0,
\]
where $h_{ij}(t)\ge0$ is the verified harm attributable to $i$'s conduct toward $j$ in period $t$, $\gamma$ converts harm to cost, and $r$ is a rehabilitation rate at which past conduct is forgiven.

\begin{proposition}[Cost Dynamics with Rehabilitation]\label{prop:relcost}
If $h_{ij}(t)\equiv h$ is constant, then $C_{ij}(t)\to\gamma h/r$. If harmful conduct ceases at time $T$, so that $h_{ij}(t)=0$ for $t\ge T$, then $C_{ij}(T+n)=(1-r)^nC_{ij}(T)\to0$. In general, $C_{ij}(t)=(1-r)^tC_{ij}(0)+\gamma\sum_{s<t}(1-r)^{t-1-s}h_{ij}(s)$ and is bounded by $\max(C_{ij}(0),\gamma\sup_sh_{ij}(s)/r)$.
\end{proposition}

\begin{proof}
The closed form follows by induction. A constant $h$ gives the geometric series $\gamma h\sum_{k<t}(1-r)^k\to\gamma h/r$, and the initial condition decays as $(1-r)^t$. Cessation of harm leaves the homogeneous recursion with solution $(1-r)^nC_{ij}(T)$. The bound follows from $\sum_{k\ge0}(1-r)^k=1/r$.
\end{proof}

The model has two implications. Cost is bounded for bounded harm, so that the principle cannot produce unlimited accumulation of penalty from a bounded stream of misconduct. And cost decays when harm ceases, so that the possibility of rehabilitation is built into the dynamics and is not an optional addition. The decay is the same arithmetic as that of the reputation score of Chapter 4, and the apparent tension with the argument against decaying scores is resolved by the difference in what is being decayed. Chapter 4 objected to decay of recognition for accomplishment, which penalizes silence. The decay here is of a penalty for harm, which releases a person from a burden when the conduct that earned it has stopped, and it is an expression of the right to finish applied to wrongdoing.

\section{The Danger of a Universal Reputation Score}

This section should explicitly return to Lanier. A universal relationship-cost ticker could become an intrusive social scoring system, recreating the very pathology the essay seeks to avoid. If every person carried a continuously updated cost, available to every counterparty, the cost would acquire the properties that Chapter 2 identified as the source of the compulsion: high salience, continuous updating, and wide use in determining opportunity. People would conduct themselves with a view to the score, the score would shape their preferences according to the dynamics of Chapter 3, and the score would be gamed in the manner of Chapter 3. The very capacity that makes the principle valuable, namely that information about conduct can be conveyed to those who deal with the agent, is the capacity that makes a social credit system possible.

Two further risks attend any such system. Concentration of power is one: a price is paid more easily by those with greater resources, so that wealthy actors may be able to absorb the costs and continue, and the instrument then deters only the poor. The remedy is to scale costs to the gain from the transaction and to reserve prohibition for conduct that no price should license, as Proposition \ref{prop:friction} indicates. Institutional incentive is the other: the bodies that assess costs must not benefit from imposing ever more of them. An assessor funded by the proceeds of the costs it imposes is in the position of the institution funded by the persistence of the problem in Proposition \ref{prop:funding}, and will find reasons to discover new offenders. The proceeds should therefore not flow to the assessing body, and the assessor should be evaluated by the accuracy of its findings on review and not by their number. The danger is not removed by good intentions in the design. It is structural. A single number summarizing a person's conduct across all domains requires a common scale of commensuration, and commensuration across domains is exactly the operation that Simmel and Espeland and Stevens \cite{espeland} identify as both powerful and distorting. The essay's recommendation is accordingly to avoid the aggregate. Relationship costs should be domain-specific, so that conduct in the supply of one kind of goods bears on dealings in that domain and not on unrelated dealings. They should be relationship-specific where possible, so that the cost is borne in the terms between particular parties. And they should be accessible only to those with a legitimate need in the domain, so that the information does not circulate beyond its purpose.

\section{Procedural Safeguards}

Six requirements follow for any system of relationship costs. Contestability requires that the affected party be informed of the findings on which a cost rests and given a genuine opportunity to challenge them before an impartial body, since findings made without a chance of rebuttal are prone to error and abuse. Proportionality requires that the cost be commensurate with the verified harm and the risk it indicates, so that minor failings do not trigger severe consequences. Evidence requires that findings rest on independently established facts of a standard appropriate to the consequence, and not on allegation, rumor, or inference from association. Privacy requires that the information be used for its stated purpose, held no longer than needed, and shared only with those with a legitimate need. Time limits require that costs expire, in accordance with the decay of Proposition \ref{prop:relcost}, unless renewed by fresh findings. And the possibility of rehabilitation requires that there be a defined path by which an agent can demonstrate changed conduct and have the cost reduced before it has run its course. Two further limits concern reach. The cost should be limited in scope, applying only where a transaction meaningfully contributes to the harm, and not to every dealing with a party who has also engaged in harmful conduct elsewhere. And indirect costs should decline with each step through the network, in accordance with Proposition \ref{prop:decay}. Taken together, these requirements are the elements of due process, and the claim is that a system of relationship costs without due process is a system of arbitrary penalty, whatever its stated purpose.

\section{Local Accountability Rather Than Universal Surveillance}

Narrowly scoped, purpose-specific accountability may be preferable to generalized social valuation. The argument is both practical and principled. Practically, narrow systems are easier to make accurate, since the relevant conduct can be defined and verified, and easier to contest, since the affected party knows what is at stake. They also contain the consequences of error, since a mistaken finding in one domain does not follow the agent into others. On principle, a narrow system respects the plurality of human life in which a person may be excellent in one respect and deficient in another, and does not require that these be reduced to a single standing. The essay's position is therefore that accountability should be built in the form of many specific, contestable, time-limited, and domain-bound practices and not a single comprehensive account, and that the difference between the two is the difference between a constitution of limited powers and a surveillance apparatus.

\part{The Architecture of a Quieter Information System}

\chapter{From Broadcasting Problems to Routing Responsibility}

\section{The Broadcast Model}

In the broadcast model, information is distributed according to reach, engagement, or visibility. A problem is published to a feed, and its prominence is determined by the attention it attracts. This arrangement has a certain democratic appeal, since it allows any observer to bring a matter before the public, and it has produced real benefits, including the exposure of wrongdoing that institutions had concealed. Its structural feature is that the selection of the recipients of a message is indifferent to their capacity to act. Delivery is governed by the properties of the message, such as its capacity to provoke reaction, and not by the properties of the recipients.

\section{The Routing Model}

In the proposed alternative, a problem is directed toward agents capable of evaluating and addressing it. The model is borrowed from the organization of any functioning institution. A hospital does not broadcast a patient's condition to all staff but routes it to the physician responsible, with escalation to specialists as the case requires. A software project does not announce each defect to all contributors but assigns it to a maintainer of the affected component. The routing model presupposes a representation of the recipients' competence and of the structure of the problem, and the quality of the system depends on the quality of that representation.

\section{Competence-Based Routing}

An information system could identify appropriate recipients according to expertise, authority, availability, and feedback capacity. The four attributes correspond to the dimensions of the readiness index of Chapter 8. Expertise concerns the capacity to diagnose and intervene. Authority concerns legitimacy. Availability concerns whether the recipient can take on the task in the time required. Feedback capacity concerns whether the recipient is in a position to verify the outcome. W. Ross Ashby's law of requisite variety \cite{ashby} supplies a general reason that routing must be matched to the diversity of problems: a regulator can control only as much variety in a system as it has variety of its own responses. A routing system with few categories cannot handle problems of many kinds, and a system that uses a finer classification requires a correspondingly large and diverse set of competent agents.

A comparison of the two models under limited attention can be stated in a simple form. There are $M$ problems, each requiring attention from at least one competent agent. There are $N$ agents, of whom a fraction $q$ are competent for any given problem. Each agent can consider $k$ problems in a period.

\begin{proposition}[Broadcast and Routed Detection]\label{prop:route}
In the broadcast model each agent considers $k$ problems drawn uniformly without replacement from the $M$ in a common feed, and a problem is addressed if at least one competent agent considers it. The probability that a given problem is addressed is
\[
P_b=1-\Bigl(1-\frac kM\Bigr)^{qN},
\]
which satisfies $1-e^{-qNk/M}\le P_b\le\min(1,qNk/M)$. In the routed model, a classifier directs each problem to a single competent agent with probability $r$ (and to no competent agent otherwise), and each competent agent receives at most $k$ problems in expectation, which requires $rM\le qNk$. Then the probability that a given problem is addressed is $P_r=r$. Consequently, whenever $r>qNk/M$ and the load condition holds, $P_r>P_b$, and if $qNk/M\to0$ with $r$ fixed, $P_b\to0$ while $P_r$ remains $r$.
\end{proposition}

\begin{proof}
Each competent agent independently considers a given problem with probability $k/M$, so the probability that none of the $qN$ competent agents considers it is $(1-k/M)^{qN}$. The lower bound follows from $1-x\le e^{-x}$ and the upper bound from the union bound. The routed model is by construction. The comparison follows from $P_b\le qNk/M<r=P_r$.
\end{proof}

The proposition is a rigorous statement of a restricted claim. It does not establish that routing is better than broadcasting in general. It identifies a regime, in which problems are numerous relative to the attention of competent agents, and in which the classifier is better than the baseline rate of chance exposure, where routing dominates. It also exhibits the complementary regime: if the classifier is poor, with $r$ small, broadcast outperforms, and the superiority of routing is lost. A mixed system in which routing is the default and broadcast the fallback for problems that the classifier cannot place is suggested by the analysis, but the optimal mixture depends on assumptions about the classifier and about the costs of misrouting that the model does not specify. The comparison of the two architectures under heterogeneous expertise and uncertain problem classification is therefore left as the major open question and the point of contact between the philosophical proposal and empirical institutional design.

\section{Escalation Conditions}

A problem should be referred beyond the local level under specified conditions. The first is inadequate competence: the agent to whom the problem is routed cannot diagnose or remedy it. The second is insufficient resources: the agent could address it in principle but lacks the means. The third is unresolved externalities: the problem affects parties outside the agent's domain who have not been consulted. The fourth is conflict of interest: the agent has a stake in the matter that compromises judgment. The fifth is failure to achieve correction: the agent has tried and the problem persists. These conditions correspond to the exceptions in the default rule of Section 8.6, and the system should make escalation a recognized and unstigmatized step and not a sign of failure.

\section{Avoiding Bureaucratic Gatekeeping}

A routing system could suppress legitimate complaints by incorrectly declaring them outside an agent's jurisdiction. The danger is real and well documented in bureaucratic experience: the complainant is passed from office to office, each of which disclaims responsibility, and the effect is the disappearance of the complaint. The danger is the shadow of the proposal, since the same classification that routes a problem to the right place can route it to no place. The mechanisms to counter it are appeal, independent review, and public-interest escalation. Appeal allows a person whose complaint was declared out of jurisdiction to have the decision reviewed. Independent review places an authority outside the routing hierarchy in a position to examine patterns of refusal. Public-interest escalation provides that a matter of broad concern that the ordinary channels have failed to address may be brought before a body with the standing to compel attention. A routing system without these mechanisms is a mechanism for burying problems, and its efficiency in doing so would be a defect.

\section{The Distinction Between a Notification and a Demand}

Not every event requires interruption. A system should distinguish archival information, optional background awareness, and actionable requests requiring timely response. Archival information is retained for later retrieval and does not announce itself. Optional background awareness is available to those who wish to consult it and does not intrude on those who do not. Actionable requests are presented in a manner that makes clear what is wanted, from whom, and by when. The ordinary notification collapses the three into a single mode of presentation, an alert, which treats all information as a demand. Restoring the distinction means that the weight of an interruption is reserved for the matters that justify it, and the next chapter develops the design principles that follow.

\chapter{Designing Against the Never-Ending Ticker}

\section{Static Records Versus Moving Displays}

The informational properties of stable task lists, searchable archives, dashboards, and continuously scrolling interfaces differ in ways that bear on their suitability for different purposes. A stable task list shows what is outstanding and allows the user to scan, compare, and plan, and its state changes only when something changes. A searchable archive supports retrieval by query and is silent unless consulted. A dashboard summarizes the state of a system through a selected set of indicators, and its usefulness depends on the choice of indicators and the frequency of update. A continuously scrolling interface displays a rotating subset of items, so that a viewer sees any given item only when it passes. The scrolling form is ill-suited to the tasks of finding, comparing, and planning, which require simultaneous access to the whole, and it is well-suited, if at all, to the task of conveying that a system is alive. It trades the informational properties of a static display for an attentional one.

\section{Event-Driven Attention}

Interfaces can remain quiet unless a relevant state change requires consideration. An event-driven design generates a signal when a monitored quantity crosses a threshold or when a condition becomes true, and generates none otherwise. The signal is thereby informative: its occurrence means that something happened. The contrast with the ticker is that the ticker's motion is uninformative since it occurs regardless. An event-driven interface restores the correspondence between the activity of the display and the activity of the system, and thus reverses the condition of Theorem \ref{thm:motion}.

\section{Actionability Thresholds}

Notifications should be evaluated partly according to whether the recipient can take a meaningful action. A notification that cannot lead to action imposes an interruption cost and yields no benefit, and the criterion for sending should compare the two. Let events arrive at rate $\lambda$ with each event carrying an expected value $v\ge0$ to the recipient of timely notice, distributed according to a continuous distribution function $F$, and let each interruption cost the recipient a fixed amount. If the recipient can tolerate only $\bar N<\lambda$ interruptions per period, the problem is to choose which events to notify.

\begin{proposition}[Threshold Notification]\label{prop:threshold}
Among all notification rules, randomized or not, that deliver at most $\bar N$ interruptions per period in expectation, the rule that notifies exactly when $v\ge\theta$, with $\theta=F^{-1}(1-\bar N/\lambda)$, maximizes the expected value delivered.
\end{proposition}

\begin{proof}
Let a rule notify an event of value $v$ with probability $\psi(v)\in[0,1]$, subject to $\lambda\int\psi\,dF\le\bar N$. The threshold rule $\psi^\ast=\mathbf 1_{\{v\ge\theta\}}$ satisfies $\lambda\int\psi^\ast dF=\lambda(1-F(\theta))=\bar N$. For all $v$, $(\psi(v)-\psi^\ast(v))(v-\theta)\le0$, since $\psi\le1=\psi^\ast$ when $v\ge\theta$ and $\psi\ge0=\psi^\ast$ when $v<\theta$. Integrating, $\int v\psi\,dF-\int v\psi^\ast dF\le\theta\bigl(\int\psi\,dF-\int\psi^\ast dF\bigr)\le0$, where the last inequality uses $\theta\ge0$ and the budget constraint.
\end{proof}

The result is the familiar bathtub principle, and its role here is to show that bounded interruption is not a limitation on the usefulness of notification but the condition for maximizing it under a bounded attention budget. The cost of sending a low-value notification is not merely its own cost; it is the displacement of the threshold. Each additional low-value interruption admitted lowers the threshold at the margin and reduces the value delivered by the whole system.

\section{Completion and Closure}

A system should represent the completion of an obligation as a meaningful terminal state and not merely replace it with another item. In a task system, completion should be visible as such: the item is marked complete, it moves to a record of completed work, and the number of outstanding items decreases. Where the system instead continually replenishes the list, so that the removal of one item is followed by the arrival of another, no state of the system corresponds to being done. Formally, closure requires an absorbing state: a configuration in which the set of outstanding obligations is empty and from which the system does not leave in the absence of new external demands. A system with automatic replenishment from a stream of generated work lacks such a state by construction, and the user is condemned to a task that cannot be finished. The point is not that the supply of work must end but that the user should be able to see, at least for a bounded unit of work such as a day or a project, that it has been completed.

\section{Finite Work and Infinite Presentation}

Return to the original two scrolling task rows. A finite list can be represented as an endless visual process. The design problem is whether this continuous motion helps the user locate information or instead implies a false inexhaustibility of work. The evidence of the display itself supports the second reading. A viewer who wished to find a particular task would find the scrolling arrangement an obstacle, since the item must be waited for. A viewer who wished to know how much work remained would find it uninformative, since a looping row has no beginning and no end. The display answers neither question and answers a third: is the system busy? The answer is always yes. The inference that the work is inexhaustible is not a stated claim but the natural reading of a process with no terminus, and the design fails the test of honesty in precisely this respect. It misrepresents the shape of the work, which is finite, by a form that suggests infinity.

\section{Design Criteria}

Five principles can be developed and justified. Quiet defaults hold that an interface should not move, sound, or interrupt unless there is a reason, so that the baseline state is rest and every departure from it carries information. Bounded interruption holds that the number of times the system may claim the user's attention in a period is capped, with the cap chosen by the user, so that the threshold of Proposition \ref{prop:threshold} is a visible control and not a hidden parameter. Explicit closure holds that obligations have terminal states that are displayed as such, so that completion is perceivable. User-controlled retrieval holds that information not needed now is available on demand and does not push itself forward, in accordance with the distinction between consulted and soliciting measurements drawn in Chapter 2. And preservation of historical accomplishment holds that completed work and earned recognition remain available without maintenance, in the form of the accumulating record of Chapter 4.

Importantly, the principles should not be extended to the proposition that every interface must be static. Some continuous displays are justified by rapidly changing conditions, such as emergency monitoring or industrial control. A display of reactor coolant temperature, an air-traffic radar picture, or an intensive-care monitor moves because the world it represents moves, and its motion is information in the sense that Theorem \ref{thm:motion} denies to the ticker. The criterion that separates the cases is the one already stated: whether the activity of the display is a function of the state of the system. Where it is, continuous motion is faithful. Where it is not, it is ornament, and it may be worse than ornament if it trades on the viewer's inference from motion to change.

\part{Philosophical Foundations and Objections}

\chapter{Agency, Finitude, and the Ethics of Scale}

\section{The Finite Agent in an Effectively Unbounded Information Environment}

An individual possesses limited time, attention, knowledge, and causal reach. These limits are not defects to be overcome by better technique but features of what it is to be an agent among other agents in a world that is larger than any of them. The days available are finite, the number of things that can be held in mind at once is small, the knowledge accumulated in a life is a minute fraction of what is known, and the portion of the world on which a person's action has any appreciable effect is narrow. A conception of ethics that ignored these limits would be a conception of ethics for a different kind of being.

Digital systems can expose that individual to far more events than can be investigated or acted upon. The exposure is effectively unbounded in the sense that no finite person can exhaust it: new material arrives faster than any person can consider it, and the supply does not diminish with consumption. The consequence is that selection is unavoidable. A person cannot choose to attend to everything, and what she attends to is determined either by her own judgment or by the mechanisms that present material to her. The ethical significance of information design lies partly in this fact. A mechanism that selects for the person, according to criteria she did not choose, is making a decision about the use of her finite life.

\section{The Difference Between Information and Responsibility}

Possessing information does not automatically establish either authority or competence. This claim, which may seem obvious, runs against a persistent assumption of the information age, that the democratization of information is thereby a democratization of capability and of standing. The assumption has a respectable history in the thought that secrecy is the instrument of domination and that disclosure is its remedy. There is truth in this, and the argument of this essay is not directed against transparency. The error is the further inference that whoever has access to information has thereby acquired a role in the matter it concerns.

The relation of information to responsibility is better understood through the concept of Hans Jonas \cite{jonas}, who argued that technological power creates novel responsibilities because the reach of action has outrun the reach of traditional ethics, and that responsibility is proportioned to power and to foresight. On this view responsibility follows from the capacity to affect and to foresee, and not from the capacity to observe. A person who observes a distant event without the power to affect it has not thereby acquired a responsibility for it, though she may have acquired a reason to care. Conversely, a person who wields power over matters she does not observe has responsibilities, which she is bound to discharge by acquiring the relevant information. The two relations run in opposite directions: information is an obligation of power, and not a source of it.

\section{Situated Knowledge}

Friedrich Hayek's discussion of the use of knowledge in society \cite{hayek} argued that the knowledge relevant to economic decisions is not given to any single mind but is dispersed among individuals in the form of particular circumstances of time and place, and that the problem of social organization is how to make use of this dispersed knowledge. Hayek's emphasis was on the price system as a device for communicating the information needed for coordination without requiring that any agent possess the whole. Elinor Ostrom's study of the governance of common-pool resources \cite{ostrom} documented the capacity of communities to devise and sustain rules for the management of shared resources, drawing on local knowledge of conditions and on monitoring and sanctions that those closest to the resource could apply. Her work on polycentric governance, in which multiple centers of authority operate at different scales with overlapping jurisdiction, suggested that effective arrangements often combine local autonomy with higher-level coordination.

The two contributions lend support to the epistemic argument for locality, and they should be used with discrimination. The epistemic advantages of local knowledge are real: those who live with a problem know things about it that outsiders do not. But the advantages do not entail that local actors are invariably more ethical or effective. Local knowledge can be parochial, local institutions can be captured by local elites, and local norms can be unjust. Ostrom herself catalogued cases of failure along with cases of success, and identified design principles that distinguished them. The lesson to be drawn is that locality is a resource for good governance and not a guarantee of it, and that the question is how institutions can be arranged so as to make use of local knowledge while correcting local failure.

\section{The Phenomenology of Perpetual Incompletion}

The experience of living within systems that continually present new demands without corresponding mechanisms of closure can be described in terms of the structure of an unfinished action. A person who has begun a task and not completed it retains a standing orientation toward it, a readiness to resume that occupies some portion of attention. This is the ordinary condition of a life with projects, and it is manageable so long as projects are finished with reasonable frequency. A system that adds items as fast as or faster than they are removed produces a permanent state of incompletion, in which there is always an unfinished matter and the orientation toward it never relaxes.

The point is phenomenological and psychological, and the psychological claims should be held with caution. There is a body of work on the effects of interrupted and unfinished tasks on cognition, and its results have been the subject of dispute and revision, so that no strong claim is made here. The narrower point is philosophical. Action is intelligible as the kind of thing that has an end, and a practice that is structurally without an end is not a practice of action in the full sense but a practice of continuation. A person embedded in such a practice may find that she cannot say what she has accomplished, because the criteria for accomplishment have been replaced by the criteria for continuing. Norbert Wiener's analysis of communication and control \cite{wiener} emphasized feedback as the condition of purposive behavior: a system steers by comparing its state with a goal and correcting the difference. A system that has no goal state toward which it steers cannot register that it has arrived, and for such a system the notion of success has no application.

\section{Freedom from Compulsory Awareness}

Attentional autonomy, the capacity to determine what one attends to, deserves consideration as an important component of practical freedom. The thought is that a person who cannot control her attention is not in full control of her life, since attention determines what she notices, what she cares about, and what she does. Analyses of contemporary data-driven economies have argued that the capture and direction of attention has become a central object of commercial design, and that the resulting practices raise questions about autonomy \cite{zuboff}. The essay takes no position on the extent to which these analyses are accurate in detail. It notes only that the question is a serious one and that the protocol can be seen as a proposal for institutional protection of attentional autonomy.

The argument should recognize competing obligations. There is a need to remain informed about matters affecting others, which arises from citizenship, from professional roles, and from ordinary decency. A person who retreated from all awareness of public affairs in the name of attentional autonomy would be failing in obligations that do not depend on her power to affect outcomes alone. The resolution is again one of allocation. There is a difference between awareness that is chosen, structured, and proportioned to one's position, and awareness that is imposed, continuous, and indiscriminate, and the freedom at issue is the freedom to prefer the first. An informed citizen reads the newspaper, studies the ballot, and attends to the affairs of the community, and does so at times and in forms of her own choosing. She is not required to carry a ticker.

\chapter{Objections to the Mind Your Own Business Protocol}

This chapter should be among the strongest in the essay, because the proposal is vulnerable to several serious objections, and the credibility of the whole depends on how they are met. Each is stated in its most forceful form, and each is answered in a way that identifies what the objection establishes and what it does not.

\section{The Complacency Objection}

Does prioritizing local problems encourage people to ignore systemic injustice? The objection is that the protocol, under cover of a plausible doctrine of competence, provides a respectable rationale for withdrawing from the struggles that most require participation. History is full of people who attended to their gardens while injustice flourished, and many of them said what the protocol says: that it was not their business. If a principle of bounded responsibility had been widely adopted in such a society, it could have supplied the moral vocabulary of acquiescence.

The reply begins by granting the possibility of misuse. Any principle of priority can be invoked as an excuse, and the protocol's exceptions are designed to limit this, since the presence of an external harm or a failure of local accountability defeats the presumption. The reply continues by distinguishing the allocation of direct intervention from the obligations of citizenship, solidarity, and institutional oversight. The protocol concerns who should do what directly. It leaves untouched the obligations that every person has as a member of a political community, which include voting, supporting accountable institutions, refusing complicity in unjust practices, and contributing to the organizations that address systemic problems. These are not interventions in the problems themselves but participation in the collective arrangements through which problems are addressed. The protocol's contention is that such participation is better conducted through organized and accountable channels than through the individual's attempt to resolve each matter herself, and the contention is compatible with a high level of political engagement. A person who has understood the protocol should be more effective in the struggle against injustice than one who has not, because she directs her efforts where they count.

\section{The Privilege Objection}

Can powerful actors define the boundaries of their own responsibility so narrowly that external harms remain invisible? The objection points to a characteristic abuse of any doctrine of limited responsibility. The corporation that disclaims responsibility for conditions in its supply chain, the official who attends only to the narrow terms of her brief, and the investor who is not concerned with the consequences of the firms he funds are all, in a sense, minding their own business. If the protocol allows each agent to determine the extent of her business, it permits the powerful to define their business so as to exclude the harms they cause.

The objection is correct about the danger and shows that jurisdiction cannot be self-assigned. The protocol's answer is that the boundaries of responsibility are not determined by the agent's own preferences but by facts that can be independently ascertained: the agent's causal contribution to the harm, her control over the relevant processes, and her capacity to foresee the consequences. An agent who causes a harm has responsibility for it, whatever she may claim about the scope of her concern. This yields a causal contribution principle: the domain of an agent's business includes whatever she causes or controls, and the claim that a matter falls outside her business is itself subject to review. The need follows for independent jurisdictional review, a body or procedure that can examine such claims and rule on them. The doctrine of responsibility without jurisdiction was introduced to describe the position of the powerless, who are burdened with demands they cannot meet, and not to excuse the powerful, who have jurisdiction over matters they prefer to ignore. The protocol constrains the second as much as it relieves the first.

\section{The Coordination Objection}

How can complex societies address problems that no individual can understand or solve independently? The objection is that many of the most important problems, such as climate change, financial stability, pandemic response, and the governance of technology, are too large and too complex to be addressed by agents attending to local matters. They require coordination across many parties, and coordination requires that many people share a picture of the problem and a commitment to act on it. A principle that directs attention to the local would seem to fragment the shared attention that coordination requires.

The response is that coordination at scale has never been achieved by universalizing individual attention, and that the successful instances have relied on modular expertise and coordinated institutions. A large engineering project is not coordinated by every participant understanding the whole. It is coordinated by interfaces, specifications, and the division of labor, so that each participant understands her part and the way it connects to others. A public health system does not rely on every citizen understanding epidemiology but on trained professionals, institutions with defined responsibilities, and public communication calibrated to what citizens need to do. The protocol is a principle for the participants in such an arrangement: to do their part well, to communicate across interfaces, and to escalate what exceeds their competence. What it opposes is the idea that coordination can be produced by saturating everyone with the totality of information, an idea whose failures Chapter 5 analyzed. The essay's thesis is not that coordination is unnecessary but that effective coordination depends upon restoring the locality of practical agency.

\section{The Emergency Objection}

What happens when distant emergencies create legitimate demands for immediate assistance? The strongest form of the objection is the classical argument advanced by Peter Singer \cite{singer}: if one can prevent something very bad from happening at small cost to oneself, one ought to do so, and the fact that the person in need is far away does not alter the obligation. A principle that gives priority to the local seems to conflict with the claim that distance is morally irrelevant to the duty of rescue.

The protocol's reply should concede that distance does not by itself diminish the moral claim of those in need, and should therefore treat the matter as one of the exceptions to the default rule, not as a refutation of the rule. The circumstances under which ordinary locality preferences should be overridden include the severity of the harm, its irreversibility, the urgency of the need, the ability of the agent to relieve it at proportionate cost, and the availability of a reliable channel through which relief can be delivered. When these conditions hold, the emergency exception applies, and the agent is bound to act. The practical question is how the agent can act effectively, and here the protocol's emphasis on routing and institutional channels is relevant: a donation to a reliable relief organization is a form of effective intervention, because the organization has the access, competence, and feedback that the donor lacks. The protocol endorses such action. What it resists is the inference from the existence of distant emergencies to a standing obligation on every individual to monitor and respond to all of them continuously.

\section{The Epistemic Objection}

Local observation can be misleading, while carefully conducted distant research can be highly reliable. The objection is that the protocol's emphasis on observability appears to favor the testimony of one's own eyes over the findings of systematic inquiry, which is the reverse of what epistemology recommends. A person in a neighborhood may judge a policy by its apparent local effects and be wrong, whereas a large study across many places may show its real effect. Proximity yields anecdote, and distance, properly handled, yields evidence.

The objection is sound, and the essay rejects any simplistic equation between physical proximity and epistemic authority. Locality in the sense used here is relational, and it includes the institutional extension of observation described in Chapter 10. A researcher who studies a distant population by well-designed methods has a form of access to it that is more reliable than that of a neighbor. The claim of the protocol is therefore not that local perception is superior to research, but that the agent who is to intervene must have some route to reliable feedback about the effects of the intervention, and that for the ordinary person that route is most often local. It is compatible with the claim that when a question concerns the effects of a policy across a population, the right source is a study and not a glance at the street. The protocol assigns the person who will act the responsibility of consulting the best available evidence, and it allocates to the researcher the task of generating it.

\section{The Measurement Objection}

Does the protocol merely replace one system of numerical optimization with another? This is arguably the deepest objection. The readiness index of Chapter 8 and the allocation model of Chapter 9 are numerical, and a reader might fear that they presage a regime in which individuals are scored and ranked by their actionable responsibility. The proposed model of actionable responsibility must not become a universal score governing the individual's worth.

The response has three parts. First, the apparatus is analytic. It is employed to understand the structure of allocation problems and to evaluate institutional designs, and the individual is not asked to compute it. Second, the protocol can be implemented without any scalar score at all. Its operative content can be expressed as a set of Boolean eligibility conditions, that an agent has access, capacity, feedback, and authorization, together with coarse categories and defeasible rules. Such a rule system assigns no number to any person and induces no ranking among persons. It says of a given problem whether a given agent is a suitable recipient, and it says nothing of the agent's merit. Third, the safeguards of Chapter 12 apply. No aggregate score should be computed, displayed, or used to determine opportunity. If an institution observed that it was computing one, the protocol would condemn the practice by its own standards. The objection is thus correct that the risk exists and its solution is institutional discipline, not mathematical form. A theory can use numbers without prescribing that anyone be numbered.

\section{The Democratic Objection}

Could limiting attention to one's immediate competence weaken democratic deliberation about collective decisions? The objection is that democracy requires an informed citizenry, and that a doctrine encouraging citizens to attend to their own affairs would erode the public engagement on which self-government depends. The citizen who has been persuaded that only the local is her business may feel that national and international questions are not hers to judge, and abdicate her role in deciding them.

The answer distinguishes competent participation in collective governance from the expectation that every citizen should continuously monitor every controversy. Democratic competence is not omniscience. It is the capacity to participate in the decisions that belong to citizens, to choose representatives, to judge the broad direction of policy, to hold officials to account, and to rely appropriately on trustworthy sources for matters beyond one's expertise. These activities are discrete and periodic, and they are better served by deliberate preparation than by continuous exposure. An election is an occasion, a hearing is an occasion, a referendum is an occasion, and the attention appropriate to each is of a different quality from that induced by a perpetual feed. There is, moreover, a case that continuous exposure harms democratic deliberation, if it produces the shallow judgment and misplaced confidence discussed in Chapter 6, and if it substitutes the performance of opinion for the exercise of choice. The protocol therefore claims to serve democratic deliberation, by freeing the attention that deliberation requires from the distractions that compete with it.

\part{Synthesis and Conclusions}

\chapter{The Difference Between Being Busy and Doing Something}

\section{Six Propositions Kept Apart}

The chapter returns to the distinctions with which the essay began and shows how they connect. Representational movement is not necessarily state change. State change is not necessarily improvement. Improvement is not necessarily attributable to the person reporting it. Awareness is not necessarily competence. Competence is not necessarily legitimate authority. And continuous participation is not necessarily evidence of responsibility.

These are logically separate propositions, and the chapter should develop them as such and not compress them into a single rhetorical opposition between busyness and action. The compression is tempting, and it is how the point is usually made, but it loses the structure. Each of the six non-implications corresponds to a specific way in which a system can mislead, and each requires a specific remedy. To make the independence precise, introduce propositional atoms $M$ (representational movement), $S$ (state change), $I$ (improvement), $T$ (attributable to the reporter), $K$ (awareness), $C$ (competence), $L$ (legitimate authority), $P$ (continuous participation), and $R$ (responsibility), and consider the six conditionals
\[
\varphi_1=M\to S,\quad \varphi_2=S\to I,\quad \varphi_3=I\to T,\quad \varphi_4=K\to C,\quad \varphi_5=C\to L,\quad \varphi_6=P\to R .
\]
The claim that none of these holds in general is the claim that each is falsifiable. The stronger claim that they are mutually independent is the following.

\begin{proposition}[Independence of the Six Inferences]\label{prop:indep}
For each $i\in\{1,\dots,6\}$ the set $\{\varphi_j:j\ne i\}$ does not entail $\varphi_i$.
\end{proposition}

\begin{proof}
For each $i$ it suffices to exhibit a valuation that falsifies $\varphi_i$ and satisfies all the others. Atoms not mentioned are assigned false. For $\varphi_1$: let $M$ be true and $S$ false; then $\varphi_2$ and $\varphi_3$ hold because $S$ and $I$ are false, and $\varphi_4$, $\varphi_5$, $\varphi_6$ hold because their antecedents are false. For $\varphi_2$: let $M$ and $S$ be true and $I$ false; then $\varphi_1$ holds, $\varphi_3$ holds since $I$ is false, and the rest hold vacuously. For $\varphi_3$: let $I$ be true and $T$ false, with $M$ and $S$ false; then $\varphi_1$ holds vacuously, $\varphi_2$ holds because its consequent $I$ is true, and the rest hold vacuously. For $\varphi_4$: let $K$ be true and $C$ false; the others hold vacuously, including $\varphi_5$ since $C$ is false. For $\varphi_5$: let $C$ be true and $L$ false, with $K$ false; $\varphi_4$ holds vacuously and the others hold vacuously. For $\varphi_6$: let $P$ be true and $R$ false; the others hold vacuously.
\end{proof}

The proposition is elementary, as all such independence results are, and its use is to block a specific fallacy of composition. A claim of the form that continuous participation shows responsibility, or that a moving display shows progress, or that a person who knows of a problem can solve it, is not a theorem of logic or of any plausible background theory, and it requires bridging premises that must be argued for in each case. The system designer who presents movement as evidence of progress is relying on an unargued bridge. The reader who accepts it is accepting the bridge without examination.

\section{The Accountability Chain}

The principal synthesis is that an effective institution must maintain a defensible relationship between the problem it represents, the people authorized to address it, the actions undertaken, and the evidence that correction occurred. Write this as a formula. An institution $\Pi$ is \emph{accountable} for problem $p$ when
\[
\mathrm{Acc}(\Pi,p)\ \equiv\ \exists a\,\exists x\,\exists e\ \bigl[\mathrm{Rep}(\Pi,p)\wedge\mathrm{Auth}(a,p)\wedge\mathrm{Did}(a,x,p)\wedge\mathrm{Ver}(e,x,p)\bigr],
\]
where $\mathrm{Rep}$ says that the institution faithfully represents the problem, $\mathrm{Auth}$ that agent $a$ is authorized to address it, $\mathrm{Did}$ that $a$ undertook action $x$ directed to it, and $\mathrm{Ver}$ that evidence $e$ independently verifies that $x$ corrected $p$. Each conjunct can fail separately, and the failures correspond to familiar pathologies. If $\mathrm{Rep}$ fails the institution misdescribes the problem, as the moving display misdescribes finite work. If $\mathrm{Auth}$ fails, action is taken by those without standing. If $\mathrm{Did}$ fails, there is representation and authority and nothing done. If $\mathrm{Ver}$ fails, something is done and no one can say whether it worked. The institution that reports activity in place of correction has typically retained $\mathrm{Did}$ and lost $\mathrm{Ver}$, and the institution that broadcasts awareness has retained $\mathrm{Rep}$ and lost the rest.

The conjunction is no more than the formalization of an ordinary standard of accountability, and its value lies in providing a checklist that can be applied to any institution, interface, or practice. To ask of a dashboard whether it is part of an accountable arrangement is to ask which of the four it supplies and which it presumes. The moving ticker supplies none of them: it does not represent the state of the problem, it does not identify those who are authorized, it does not report what was done, and it does not verify the result. It supplies the impression of the fourth by the persistence of its motion.

\section{Coherence of the Argument}

The structure of the whole argument can now be summarized. The first part established that representation can move without the world changing, and that the motion therefore carries no information about it (Theorem \ref{thm:motion}). The second part showed that measurement enters the loop it measures, that a visible weight on a proxy displaces value toward the measured (Propositions \ref{prop:monotone}, \ref{prop:displace}), that temporal architecture can create perpetual obligation by arithmetic (Propositions \ref{prop:maintain}, \ref{prop:inactive}), and that attention is a finite resource whose saturation degrades the processing of what matters (Proposition \ref{prop:saturation}). The third part argued that the globalization of concern produces responsibility without jurisdiction and that bounded problems are the setting in which agents can learn. The fourth proposed the protocol as an allocation rule and proved the properties of its readiness index (Proposition \ref{prop:readiness}) and the structure of the optimal allocation (Propositions \ref{prop:equimarg}, \ref{prop:local}). The fifth considered institutions and showed that funding by the persistence of a problem rewards its preservation (Proposition \ref{prop:funding}). The sixth outlined an architecture of routing and quiet (Propositions \ref{prop:route}, \ref{prop:threshold}). The seventh defended the proposal against its objections. Each step is conditional on explicit assumptions, and the argument as a whole is a cumulative case and not a deduction: the premises at each stage make the conclusion more plausible without establishing it from first principles.

\chapter{The Right to Finish}

\section{A Counterprinciple}

The concluding chapter introduces the right to finish as a counterprinciple to the never-ending ticker. A task can be completed. A repair can be verified. A responsibility can be discharged. A record can preserve the accomplishment without requiring its continual reenactment. These four sentences state the content of the right, and each corresponds to a design requirement that earlier chapters developed. The completion of a task requires a terminal state (Chapter 14). The verification of a repair requires a route to feedback (Chapter 7). The discharge of a responsibility requires a boundary that defines when it has been met (Chapters 8 and 9). And the preservation of the accomplishment requires an accumulating record that is invariant under inactivity (Chapter 4).

The ticker, by contrast, is structurally indifferent to completion. Its presentation continues regardless of whether the underlying activity has reached a satisfactory endpoint. This is not a malfunction but a consequence of its construction: a display whose motion is generated by a clock has no access to the question whether anything has been finished, and a system built to maximize engagement has no interest in the answer. The right to finish is accordingly a right against a kind of design, and the claim of the essay is that this design is neither inevitable nor neutral.

\section{The Elements of the Right}

The right can be stated in terms of the three conditions that a system must satisfy to honor it. The first is the existence of closure: for a bounded unit of work, there is a state of the system that represents the work as complete, and the system can reach it. The second is the persistence of the record: the fact of completion, and the recognition earned by it, remain in force without further activity from the person, which is inactivity invariance in the sense of Proposition \ref{prop:inactive}. The third is the requirement of cause for reopening: a completed obligation is returned to the open state only by an event that supplies new evidence that it was not discharged or that circumstances have changed in a relevant way, and not by the mere passage of time or by the system's need for further engagement.

The third condition can be stated as a constraint on the transitions of a state machine. Let the state of an obligation be either $\mathsf{Open}$ or $\mathsf{Done}$. The system admits the transition $\mathsf{Open}\to\mathsf{Done}$ upon verified discharge and the transition $\mathsf{Done}\to\mathsf{Open}$ only upon an event $\varepsilon$ carrying evidence that satisfies a stated reopening criterion. The system does not admit the transition on a timer, on a change in the user's score, or on the arrival of a new promotional demand. The rule is that nothing that is complete is reopened without reason, and the reason must be one that the person can inspect and contest. It is a modest requirement, and its absence from much contemporary software is a measure of how far the design of attention-seeking interfaces has departed from ordinary standards of respect for completed work.

\section{Continuity and Its Kinds}

The essay should close by distinguishing the legitimate continuity of collective life from the unnecessary continuity of individual demands. Collective life does not end. Institutions persist, generations succeed one another, the work of maintaining a society is never finally done, and no one proposes a right to finish that would apply to a civilization. There are standing duties, and they are standing because the need they answer is permanent: the care of children, the maintenance of infrastructure, the administration of justice. This continuity is borne by institutions through division of labor, and the individuals who serve them are able to finish their shifts, hand over their tasks, and retire. The continuity of the whole is made compatible with the closure of the parts, and this compatibility is one of the principal achievements of organized social life.

The continuity that the ticker imposes is of a different kind. It places on the individual a demand that is continuous in her own experience, with no handover and no end of shift, and it derives its force from the individual's inability to see the whole of which she is supposedly a part. The distinction is between a continuity that is carried by structure and a continuity that is carried by the person. The first is the condition of a functioning society. The second is the condition of a person who is permanently on call to a world that cannot be on call to her.

\section{Tensions That Remain}

The argument will be strongest if it refuses to resolve its own contradictions prematurely, and three tensions are therefore preserved. The first is the accounting paradox. Lanier wants contributions to be recognized and compensated, yet the resulting accounting can become a source of compulsive optimization. Eliminating measurement would also eliminate some of the accountability that the system is intended to provide. The essay's resources for this tension are the separation of recognition from salience and the inactivity-invariant record, and neither dissolves it. A recorded contribution is still a measured contribution, and a person who knows she is recorded may behave differently.

The second is the locality paradox. Local observation and intervention often improve practical competence, but many consequential harms are nonlocal, distributed, and invisible to their victims. A defensible locality principle therefore requires institutions capable of extending observation and responsibility beyond immediate experience, and such institutions are themselves remote from those they serve and subject to the failures of remote institutions. The essay's resources are the mechanisms of independent verification, appeal, and escalation, and they are imperfect.

The third is the protocol paradox. A system designed to protect people from perpetual optimization could itself become another optimization system. If every individual receives a numerical score for actionable responsibility, the Mind Your Own Business Protocol has merely invented a new ticker. The protocol should not be designed to maximize an individual's measured helpfulness. It should establish boundaries within which people can make judgments, discharge responsibilities, refuse inappropriate demands, and recognize satisfactory completion. This distinction separates an institutional constitution from a behavioral optimization algorithm. A constitution limits what may be done and leaves the use of liberty to those it protects, while an optimization algorithm specifies what is to be maximized and measures the distance to it. The protocol belongs to the first kind or it fails.

\section{Conclusion}

The objective is not a world without ongoing problems. It is a world in which problems are directed toward those capable of addressing them, interventions can be evaluated, and completed obligations are allowed to remain complete. The opening example, two rows of tasks drifting across a screen, was chosen to be modest. It reveals, when examined, an interface that keeps moving after its information has stopped changing, and the examination has led through the mathematics of representation, the psychology of metrics, the economics of attention, the ethics of responsibility, and the theory of institutions to a general account of the relation between measurement, attention, responsibility, and the organization of modern life. Visibility becomes an illegitimate demand upon finite agency when it is detached from competence, feedback, and authority, and it can instead support effective correction when it is routed to those who can act and tied to evidence that the action worked.

The essay closes with a formulation of its position.

\begin{quote}
A civilization does not become more responsible by ensuring that everyone is perpetually aware of everything. It becomes more responsible by ensuring that consequential problems are observable to those competent to address them, that interventions remain accountable to their outcomes, and that people are permitted to finish what they have undertaken.
\end{quote}

\begin{thebibliography}{99}

\bibitem{ashby} W. R. Ashby, \textit{An Introduction to Cybernetics}, Chapman and Hall, London, 1956.

\bibitem{campbell} D. T. Campbell, Assessing the Impact of Planned Social Change, \textit{Evaluation and Program Planning} 2 (1979), no.\ 1, 67--90.

\bibitem{ellul} J. Ellul, \textit{The Technological Society}, trans.\ J. Wilkinson, Alfred A. Knopf, New York, 1964. Originally published as \textit{La Technique ou l'Enjeu du Si\`ecle}, 1954.

\bibitem{espeland} W. N. Espeland and M. L. Stevens, A Sociology of Quantification, \textit{European Journal of Sociology} 49 (2008), no.\ 3, 401--436.

\bibitem{ferster} C. B. Ferster and B. F. Skinner, \textit{Schedules of Reinforcement}, Appleton-Century-Crofts, New York, 1957.

\bibitem{goodhart} C. A. E. Goodhart, Problems of Monetary Management: The U.K. Experience, in \textit{Papers in Monetary Economics, Volume I}, Reserve Bank of Australia, Sydney, 1975.

\bibitem{hayek} F. A. Hayek, The Use of Knowledge in Society, \textit{American Economic Review} 35 (1945), no.\ 4, 519--530.

\bibitem{holmstrom} B. Holmstr\"om and P. Milgrom, Multitask Principal-Agent Analyses: Incentive Contracts, Asset Ownership, and Job Design, \textit{Journal of Law, Economics, and Organization} 7 (1991), special issue, 24--52.

\bibitem{jonas} H. Jonas, \textit{The Imperative of Responsibility: In Search of an Ethics for the Technological Age}, University of Chicago Press, Chicago, 1984. Originally published as \textit{Das Prinzip Verantwortung}, 1979.

\bibitem{kerr} S. Kerr, On the Folly of Rewarding A, While Hoping for B, \textit{Academy of Management Journal} 18 (1975), no.\ 4, 769--783.

\bibitem{kleinrock} L. Kleinrock, \textit{Queueing Systems, Volume I: Theory}, Wiley, New York, 1975.

\bibitem{lanier} J. Lanier, \textit{Who Owns the Future?}, Simon and Schuster, New York, 2013.

\bibitem{little} J. D. C. Little, A Proof for the Queuing Formula $L=\lambda W$, \textit{Operations Research} 9 (1961), no.\ 3, 383--387.

\bibitem{marchsimon} J. G. March and H. A. Simon, \textit{Organizations}, Wiley, New York, 1958.

\bibitem{ostrom} E. Ostrom, \textit{Governing the Commons: The Evolution of Institutions for Collective Action}, Cambridge University Press, Cambridge, 1990.

\bibitem{pigou} A. C. Pigou, \textit{The Economics of Welfare}, Macmillan, London, 1920.

\bibitem{pius} Pius XI, \textit{Quadragesimo Anno}, encyclical letter on the reconstruction of the social order, 1931.

\bibitem{power} M. Power, \textit{The Audit Society: Rituals of Verification}, Oxford University Press, Oxford, 1997.

\bibitem{simmel} G. Simmel, \textit{The Philosophy of Money}, trans.\ T. Bottomore and D. Frisby, Routledge, London, 1978. Originally published 1900.

\bibitem{simon56} H. A. Simon, Rational Choice and the Structure of the Environment, \textit{Psychological Review} 63 (1956), no.\ 2, 129--138.

\bibitem{simon71} H. A. Simon, Designing Organizations for an Information-Rich World, in M. Greenberger (ed.), \textit{Computers, Communications, and the Public Interest}, Johns Hopkins Press, Baltimore, 1971, 37--72.

\bibitem{singer} P. Singer, Famine, Affluence, and Morality, \textit{Philosophy and Public Affairs} 1 (1972), no.\ 3, 229--243.

\bibitem{strathern} M. Strathern, `Improving Ratings': Audit in the British University System, \textit{European Review} 5 (1997), no.\ 3, 305--321.

\bibitem{tue} Treaty on European Union, Article 5, consolidated version, \textit{Official Journal of the European Union} C 202, 7 June 2016.

\bibitem{weber} M. Weber, \textit{Economy and Society: An Outline of Interpretive Sociology}, ed.\ G. Roth and C. Wittich, University of California Press, Berkeley, 1978.

\bibitem{wiener} N. Wiener, \textit{The Human Use of Human Beings: Cybernetics and Society}, Houghton Mifflin, Boston, 1950.

\bibitem{zuboff} S. Zuboff, \textit{The Age of Surveillance Capitalism: The Fight for a Human Future at the New Frontier of Power}, PublicAffairs, New York, 2019.

\end{thebibliography}

\end{document}
